% Modernised reading — fonds Alexandre Grothendieck, shelfmark 14
% (pages 1-131, the whole folder).
%
% This is an INTERPRETATION, not a transcription. It re-reads the pages in
% today's notation and vocabulary, states the hypotheses the manuscript leaves
% implicit, and drops the critical apparatus so that the mathematics reads
% straight through. Everything the brackets and underlines carried moves into a
% footnote. It is held to being mathematically correct as it stands; where the
% manuscript is loose or mistaken, the true statement is what appears here and
% the footnote says what the page has. In any dispute of substance the
% transcription governs: whoever cites Grothendieck must cite batch-01.fr.tex
% (pp. 1-20) ... batch-07.fr.tex (pp. 121-131).
%
% Pass: Opus 5.5 (claude-opus-5-5), 2026-09-23 — first pass, unchecked against the pages by a human.
% The folder was read whole, all seven batches together; this is one document
% for the shelfmark. It was made on Opus 5.5 and has not been measured against
% the reference reading of folder 115 (made on Opus 5).
%
% Scope. Of the 131 leaves, about ninety carry his hand, plus seven covers
% whose words become headings here (pp. 1, 16, 39, 47, 90, 92, 101). The rest
% are blank, or typescripts used as paper (an English typescript on the Brauer
% group, a Bourbaki draft, typescripts on Hopf algebras, topos cohomology,
% sites, complete local rings, an EGA leaf, SGA 3 leaves), and p. 42, a faint
% duplicate of the letter's first leaf. None of them is mentioned. The letter
% to Serre of 27.8.65 (pp. 43-46) is condensed in the transcription because it
% is published; it is read here from what the transcription has.
%
% Spine. No single construction: a web of conjectures on algebraic cycles
% around one question — do cycles obey the laws that Lefschetz's theorems give
% the topology? — with the one degree where the answer is a theorem (degree 1,
% abelian varieties) worked out in full.
%   pp. 1-14    effectivity structures on motives, read on the motivic Galois
%               group: effective cocharacters, cones, the torus case.
%   pp. 16-20   « Équivalence d'Albanese et de Picard »: one computation whose
%               sign depends on a choice (« des signes différents !! »).
%   pp. 23-38   « Théorie de l'équivalence d'Albanese des cycles », §§1-3:
%               the ring A_alb, families parametrised by abelian varieties,
%               l-adic homomorphisms (Th. 3.1 « dém. pas faite »).
%   pp. 39-41   Tate's trick: Tate classes generated by divisor classes.
%   pp. 43-46   the letter to Serre: conjectures A, A', A'', B.
%   pp. 47-70   « Sorites sur VA, Picard, Albanese », §§1-2: correspondences
%               act on Albanese and Picard varieties; hard Lefschetz for H^1
%               with algebraic inverse (Cor. 2.6), and its converse (2.7).
%   pp. 72-73   loose sheets: Picard, Albanese and integral cohomology.
%   pp. 75-88   run « I »: an axiomatic of abelian varieties up to isogeny.
%   pp. 90-91   the table of the rings of cycle classes.
%   pp. 92-99   Mumford-Nakai, ample, effective and nef cones.
%   pp. 101-122 Lefschetz conjectures: surfaces, local, global for Chow
%               groups, checked for n = 2, 3, 4; bounded families of cycles.
%   pp. 125-127 a chain of statements on cycles algebraically equivalent to 0.
%   p. 130      Chow groups of the generic member of a linear system.
%
% Moutures kept distinct. The two numbered runs (pp. 23-38 and pp. 47-69)
% reuse the same numbers (Th. 1.6 p. 29 and Th. 1.6 p. 55; Cor. 2.2 p. 33 and
% Cor. 2.2 p. 61; Prop. 2.7 p. 34 and Prop. 2.7 p. 67): every citation here
% carries its page. The Lefschetz-type conjectures come three times: in the
% letter (A, A', A'', B, for algebraic equivalence tensored with Q), locally
% (pp. 105-108) and globally for Chow groups (pp. 110-115).
%
% What only a folder-wide reading can say (ours, and said so in the body):
%   - the blank reference « [ , 1.2.] » of Th. 1.4 (p. 28) is filled by
%     Cor. 1.2 of p. 51 (a correspondence kills Albanese kernels).
%   - Rem. 1.5 (p. 28), the NB of p. 29, the hypothesis 2) of p. 117, the
%     chain of p. 125, and the provisos « si équiv. ... = équiv. d'Albanese »
%     of pp. 114-115 are one question: does Albanese equivalence of 0-cycles
%     coincide with rational equivalence? Mumford (1969): no, over C, as soon
%     as p_g > 0. The p. 114 « ok » for 0-cycles hold only under it.
%   - the cycle U of p. 31 (« on ne sait pas si on peut trouver U ») is
%     Question 1) of p. 70.
%   - Th. 3.1 (p. 37) is what the independence left unchecked on p. 36 needs.
%   - the « ? » marked « question numérique » on pp. 114-115 are exactly
%     where Griffiths (1969) and Clemens (1983) refute the statements; the
%     same examples refute A and A' of the letter.
%
% NOTATION fixed for the folder:
%   - k algebraically closed, X smooth projective connected of dimension n,
%     Y a smooth hyperplane section; l a prime =/= char k; H^i = l-adic.
%   - CH^i(X) = cycles mod rational equivalence (his A^*, script A_r, A; his
%     « équivalence linéaire »); A^i_~(X) = Z^i(X)/~ for ~ = alb, alg, tau,
%     hom, num (his A'^*, script A_a, A_tau, A_num...); lower index = dimension.
%     Subgroups of ~'-trivial classes: CH^i(X)_alg (his J^i, P^i, A^a),
%     CH^i(X)_tau, A^i_alb(X)_alg (his P^i of pp. 33-36).
%   - Alb_X, alb_X (his A_X, S_X); Pic^0_X = (Pic^0_{X/k})_red (his B_X =
%     Pic^00); Alb_X = dual of Pic^0_X; A-hat the dual abelian variety (his A^*,
%     as in folder 12); phi_X = alb_X^* : H^1(Alb_X) -> H^1(X); psi_X :
%     H^{2n-1}(X)(n-1) -> H^1(Pic^0_X) by the Poincaré class.
%   - correspondences: for Z on X x Y, Z_* goes from Y to X on Chow groups and
%     on cohomology, {}^tZ_* from X to Y. The page writes a(gamma(Z)) in both
%     directions (p. 55 X -> Y, p. 58 Y -> X); normalised, footnoted once.
%   - pp. 1-14: Q(-1) of weight 2 (his Lambda(-1)), eps its character, y the
%     weight cocharacter (as in folder 12; his i on p. 12), i_1 the Hodge
%     cocharacter (<i_1, eps> = 1).
%   - cones (pp. 95-99): Eff, Pos, Amp, Nef^i = (Eff_i)^dual (his P^+, Q^+, R,
%     P_i^{+o}).
%
% Substantive departures, each footnoted where it occurs:
%   p. 2    « réciproquement » (Q vs C) holds for Aut(C)-stable structures;
%           the finiteness axiom read inside one finitely generated
%           subcategory (false for all motives).
%   p. 4    E(n) -> E(-n), as on p. 2.
%   p. 6    G connected supplied; pi_1 acts on Gamma' (page: Gamma).
%   p. 8    the unspecified saturation condition is automatic for closed
%           rational polyhedral W-invariant cones (ours).
%   p. 10   I(G) -> Eff(G) factors through Galois orbits (ours).
%   p. 12   dim S = 1 proved under polarisation (ours); his attempt is struck.
%   p. 14   the « fausse » of the very uncertain last paragraph: finiteness
%           holds iff the conjugates of i_1 span X_*(T)_Q (ours).
%   pp. 18-20 K read as the degree; p_*(Theta^2) = -2[0 x P] on Y; c = A x (a);
%           Delta^2 = 0 rationally.
%   p. 25   z + z' (announced z - z').
%   pp. 31-32 « U(a) = n t » holds up to Albanese equivalence, which suffices.
%   p. 33   Cor. 2.2: « diviseur » -> cycle, on A x X; Cor. 2.4 proof supplied.
%   p. 40   T_K(A) = T_l(A) (x) K.
%   p. 50   case II: psi_X (page: phi_X), twist (n-1) supplied.
%   p. 51   « projectives » (struck) restored; p. 53 Lemme 1.4 « diviseur » ->
%           0-cycle; p. 58 phi_X^{(1)} -> phi_Y.
%   p. 63   best denominator: exponent of Ker u (page: card, illegible).
%   p. 69   orientation of phi_X chosen.
%   pp. 79-84 u : A -> A' throughout (p. 79 opens with A' -> A); C^{n'-n}(A');
%           H^1(A') (page H^1(A'^*)); Lambda^{2n-i} (page 2n-1).
%   p. 86   pairing to Q_l(-1) (page Q_l(1)).
%   p. 95   x effective iff h^0(nx) >= 1 for some n (page: >= 2 for n large).
%   p. 107  N = 4: Pic(X) -> Pic(Y) injective, not bijective (3-dim node).
%   p. 110  (1), 2i = n-1: A^i_alg bijective is false (quadric in P^3); read
%           as the subgroup CH^i_alg.
%   p. 112  dim E = n+1 (page 2n+1); primitive parts are kernels of
%           L^{n-2i+1}, L^{n-2i+2} (page: L^{n-2i}, L^{n-2i+1}).
%   p. 114  n = 2 check of (3), second family: mis-indexed; restored.
%   p. 120  (ii) Z in M -> M'; p. 122 Riemann-Roch sign (-1)^{i-1}.
%   p. 125  exponents read tau = alg-trivial for 0-cycles.
%   p. 130  y' -> y_y.
%
% Announced and never established: items 1°-3° of p. 2 (not in the folder);
% the saturation condition of p. 8; the definition of « équivalence de
% Picard » (pp. 16, 91); the moving argument of Prop. 1.2 (iii) (« à
% vérifier », p. 27); the independence of Z in the duality of p. 36; Th. 3.1
% (« dém. pas faite ») and its converse, and the §3 proof, which stops on
% p. 38; Lemme 1.5 (« nous nous dispensons ») and the verification of
% Prop. 2.3 (p. 61); the denominators of Prop. 2.4; the sign questions of
% pp. 61 and 77; the questions of p. 70; problem (4) and the « Pb » of
% pp. 106-108; the proof of Th. 4 (p. 113); a primitive part of CH^* (p.
% 116); the implication (ii) => (iii) of p. 119, broken off.

\documentclass[11pt,a4paper]{article}
% Le préambule commun à toutes les transcriptions du fonds Grothendieck.
%
% Il tient en une idée : **l'incertitude se compose, elle ne se résout pas.**
% Une transcription qui lisse un mot douteux a détruit la seule chose pour
% laquelle on la faisait — un lecteur doit pouvoir savoir, à chaque endroit,
% ce qui était sur le feuillet et ce qui a été ajouté. Les six macros qui
% suivent sont donc l'appareil critique, et non de la décoration.
%
% Les mêmes commandes sont reconnues par `scripts/render.mjs`, qui produit la
% vue de lecture du site. Une macro ajoutée ici doit l'être aussi là-bas,
% faute de quoi le rendu échouera bruyamment — ce qui est voulu : un rendu
% silencieusement faux est pire qu'un rendu absent.

\ProvidesPackage{grothendieck}[2026/08/08 Transcriptions du fonds Grothendieck]

\usepackage{amsmath,amssymb,amsthm}
\usepackage{mathrsfs}
\usepackage{stmaryrd}
% Les diagrammes commutatifs sont partout dans ce fonds : c'est la langue
% même de la géométrie algébrique de Grothendieck, et une transcription qui
% les rendrait en prose ne transcrirait rien.
\usepackage{tikz-cd}
% Français par défaut — c'est la langue des feuillets — mais l'anglais est
% chargé aussi : la traduction et le résumé sont des documents anglais, et la
% typographie française (espaces avant les deux-points) y serait une faute.
% Ces documents basculent par \selectlanguage{english} en tête de corps.
\usepackage[english,french]{babel}
\usepackage{fontspec}
\usepackage{geometry}
\geometry{a4paper,margin=25mm,marginparwidth=28mm}
\usepackage{marginnote}
\usepackage{xcolor}
% ulem, not soul. `soul`'s \st and \ul are built on a character-by-character
% scan that XeLaTeX's font machinery breaks: the document fails with an
% undefined control sequence rather than degrading. ulem does the same two jobs
% and is Unicode-engine safe. `normalem` stops it hijacking \emph, which the
% transcriptions use for Grothendieck's own emphasis.
\usepackage[normalem]{ulem}
\usepackage{hyperref}
\hypersetup{colorlinks=true,linkcolor=black,urlcolor=black,pdfborder={0 0 0}}

% --- Métadonnées du lot -----------------------------------------------------
% Reprises telles quelles par le script de rendu, qui les affiche en tête de
% la vue de lecture. Elles fixent ce que le fichier prétend couvrir : sans
% elles, une transcription flotte sans point d'attache dans un fonds de seize
% mille feuillets.

\newcommand{\folder}[1]{\gdef\grothFolder{#1}}
\newcommand{\batch}[1]{\gdef\grothBatch{#1}}
\newcommand{\leaves}[2]{\gdef\grothFirst{#1}\gdef\grothLast{#2}}
\newcommand{\foldertitle}[1]{\gdef\grothFoldertitle{#1}}
\gdef\grothFolder{?}\gdef\grothBatch{?}\gdef\grothFirst{?}\gdef\grothLast{?}\gdef\grothFoldertitle{}

% --- Le résumé --------------------------------------------------------------
%
% Il ouvre la lecture modernisée, sous le titre « Résumé », et il est composé
% en petit. Ce n'est pas un ornement : c'est le seul passage du document qui
% ne soit pas des mathématiques, et un lecteur doit pouvoir le reconnaître
% avant de l'avoir lu — puis le sauter s'il connaît déjà le sujet, ou s'y
% arrêter s'il ne le connaît pas. Le filet qui le suit marque la reprise du
% corps.

\newenvironment{resume}
  {\small\setlength{\parskip}{0.45em}}
  {\par\medskip\hrule\medskip}

% --- Mots-clés ---------------------------------------------------------------
%
% En anglais, en clôture du résumé de la lecture modernisée : le vocabulaire
% moderne sous lequel chercher ce que les feuillets construisent. Le manifeste
% du site (`npm run manifest`) les extrait pour étiqueter la cote entière —
% cette ligne est la source unique des tags, il n'y a pas de fichier à côté.
\newcommand{\keywords}[1]{\par\smallskip\noindent
  {\footnotesize\textsc{Keywords} --- \textit{#1}}\par}

% --- L'appareil critique ----------------------------------------------------

% Le numéro de feuillet, en marge. C'est l'ancre qui permet de régler un
% désaccord sur une formule en regardant *un* feuillet plutôt qu'une chemise
% de six cents. Le site s'en sert aussi pour tourner les pages du fac-similé
% quand on fait défiler la transcription.
\newcommand{\leaf}[1]{%
  \marginnote{\footnotesize\color{gray!70}\textsf{#1}}%
  \label{leaf:#1}\ignorespaces}

% Illisible : on ne devine pas. Un mot inventé qui se lit comme les autres est
% le pire résultat possible.
\newcommand{\ill}{\textcolor{red!60!black}{[\,\ldots\,]}}

% Lecture douteuse : le mot est donné, et signalé comme incertain.
\newcommand{\uncertain}[1]{\uline{#1}}

% Ajout de l'éditeur — une lettre manquante, un mot sous-entendu.
\newcommand{\add}[1]{\textcolor{blue!50!black}{[#1]}}

% Biffé par Grothendieck lui-même. Ce qu'il a rayé fait partie du document :
% une rature dit souvent où la pensée a changé de direction.
\newcommand{\struck}[1]{\sout{#1}}

% Note du transcripteur, sur l'état matériel du feuillet.
\newcommand{\note}[1]{\par\smallskip{\small\color{gray!60!black}\textit{[#1]}}\par\smallskip}

% Note marginale de Grothendieck, distincte du corps du texte.
\newcommand{\marginal}[1]{\par\smallskip\noindent\hspace{1em}%
  {\small\color{gray!60!black}#1}\par\smallskip}

% --- Titre ------------------------------------------------------------------

\AtBeginDocument{%
  \begin{center}
    {\large\bfseries Fonds Alexandre Grothendieck --- cote n\textsuperscript{o}~\grothFolder}\\[2pt]
    {\itshape\grothFoldertitle}\\[4pt]
    {\small feuillets \grothFirst--\grothLast\ (lot \grothBatch)}\\[2pt]
    {\footnotesize\color{gray!60!black}Transcription établie d'après le fac-similé numérisé
     par l'Université de Montpellier.\\
     Les lectures incertaines sont soulignées, les illisibles notées [\,\ldots\,],
     les ajouts entre crochets.}
  \end{center}
  \vspace{1em}}

\endinput


\folder{14}
\pages{1}{131}
\dating{1965-[vers 1971]}
\watermark{Édition de démonstration\\interprétation personnelle de l'œuvre}
\foldertitle{Théorie des cycles algébriques : conjectures de Hodge, Tate, Lefschetz, Weil : notes manuscrites (s.d.), lettre (1965) — lecture modernisée du dossier entier}

\begin{document}

\section*{Résumé}

\begin{resume}
Ces feuillets travaillent sur les \emph{cycles algébriques} : les
sous-variétés d'une variété algébrique, comptées avec des multiplicités
entières, et ce qu'on peut dire d'elles quand on les déplace.

Voici le problème. Sur une courbe, un cycle de dimension zéro --- une somme
formelle de points --- a un degré et, si l'on additionne ses points dans la
jacobienne, une « somme » ; le théorème d'Abel dit que ces deux invariants
décident exactement quand le cycle est le diviseur d'une fonction. On voudrait
la même chose en toute dimension : une liste d'invariants qui décide quand un
cycle peut être déformé en un autre, et des lois qui relient les cycles de
dimensions différentes. Plusieurs manières de « déformer » coexistent :
l'équivalence rationnelle, la plus fine ; l'équivalence algébrique ;
l'équivalence numérique, qui ne retient que les nombres d'intersection. Entre
les deux premières, ces pages en intercalent une autre, l'\emph{équivalence
d'Albanese} : on déclare nul tout cycle obtenu en transportant, le long d'une
famille algébrique, une somme de points dont la « somme » est nulle.

L'image qui porte le dossier est celle d'une tranche. Une section hyperplane
d'une variété --- une tranche plane d'une surface, par exemple --- voit,
d'après les théorèmes de Lefschetz, presque toute la forme de la variété : les
trous de petite dimension se retrouvent dans la tranche, et ceux de grande
dimension s'obtiennent en épaississant ceux de la tranche. Grothendieck
demande si les cycles obéissent aux mêmes lois que la topologie. Dans une
lettre à Serre du 27 août 1965, qui est au milieu du dossier, il énonce ces
lois en deux conjectures, A et B, qui sont l'origine de ce qu'on appelle
aujourd'hui les \emph{conjectures standard} ; il explique qu'elles donneraient
les conjectures de Weil et une théorie des motifs, et qu'il lui manque
toujours le pas décisif. Le reste du dossier en décline des variantes : sur un
anneau local, pour les groupes de cycles à équivalence rationnelle près, pour
les variétés abéliennes, et dans un tableau qui range toutes les équivalences
en une seule figure.

Plusieurs choses y sont démontrées, et elles concernent toutes le degré le plus
bas, celui où les cycles se lisent sur deux variétés abéliennes attachées à la
variété, sa variété d'Albanese et sa variété de Picard. Toute correspondance
entre deux variétés induit un homomorphisme entre leurs variétés d'Albanese,
qui ne dépend que de sa classe numérique ; et, pour la cohomologie de degré 1,
l'inverse de l'isomorphisme de Lefschetz est donné par une correspondance
algébrique --- un cas démontré de ce qu'on appelle aujourd'hui la conjecture
standard de type Lefschetz. S'y ajoutent une étude des cônes de classes amples
et effectives autour du critère numérique d'amplitude de Nakai, et, en
ouverture, des pages sur la façon dont un groupe algébrique pourrait dire
lesquels des motifs sont « effectifs ».

Ce qui frappe, à la lecture, est la manière. Les énoncés sont posés d'abord,
en série, avec leurs numéros ; puis essayés sur les plus petits cas ---
dimension 2, 3, 4 --- et chaque ligne reçoit en marge « ok », « ? » ou
« question numérique ». Un théorème porte « dém.\ pas faite » ; un calcul
s'achève sur « des signes différents !! ». La lettre à Serre se termine en
l'invitant à réfléchir au problème, que Grothendieck dit avoir attaqué le jour
même. On voit une théorie éprouvée par ses conséquences avant que ses
fondements existent, et les points d'interrogation mis exactement là où elle
cédera.

Car une part des questions a reçu réponse depuis, et souvent par la négative.
Mumford a montré en 1969 que, sur une surface complexe qui porte une 2-forme
holomorphe non nulle, l'équivalence d'Albanese des points est beaucoup plus
grossière que l'équivalence rationnelle ; Griffiths, la même année, que
l'équivalence algébrique diffère de l'équivalence homologique même à torsion
près. Les conjectures de 1965, qui portent sur l'équivalence algébrique, sont
donc en défaut sous cette forme ; reformulées pour l'équivalence homologique,
elles sont restées ouvertes. Les conjectures de Weil ont été démontrées par
Deligne en 1974, par une autre voie.

Les noms sous lesquels chercher la suite : conjectures standard de
Grothendieck, noyau d'Albanese et conjecture de Bloch, application
d'Abel--Jacobi et groupe de Griffiths, conjecture de Tate pour les variétés
abéliennes, critère de Nakai--Moishezon et cônes nef, théorèmes de Lefschetz
locaux.

\keywords{effective cocharacter, cone of effective characters, conjugacy classes of cocharacters, CM torus, Serre group, Albanese equivalence, Albanese kernel, Bloch's conjecture, infinite-dimensionality of the Chow group of zero-cycles, exterior product of zero-cycles, Albanese variety, Picard variety, divisorial correspondence, Poincaré bundle, duality of abelian varieties, algebraic equivalence, tau-equivalence, Griffiths group, Abel-Jacobi map, intermediate Jacobian, algebraic representative, Tate conjecture for divisors, Tate isogeny theorem, products of elliptic curves, Grothendieck-Serre correspondence, standard conjectures on algebraic cycles, polarization of the Picard variety, Weil pairing, fundamental class of an abelian subvariety, topological filtration on K-theory, Nakai-Moishezon criterion, Kleiman's ampleness criterion, nef cone, pseudo-effective cone, ample cone, strictly nef divisor, nef classes of higher codimension, Hodge index theorem, local Lefschetz theorems, Grothendieck-Lefschetz theorem, Samuel's conjecture, divisor class group of a local ring, affine Lefschetz theorem, vanishing cycles, affine cone, bounded families, Riemann-Roch without denominators, blow-up formula}
\end{resume}

\pagerange{1}{131}
\section*{Le fil du dossier, et les conventions}

\subsection*{Les stations}

Le dossier ne construit pas un objet unique. Il tisse, autour d'une seule
question --- les cycles obéissent-ils aux lois que les théorèmes de Lefschetz
donnent à la topologie ? ---, un réseau de conjectures, et il traite à fond le
seul degré où la réponse est un théorème : le degré 1, celui des variétés
abéliennes.

\begin{enumerate}
  \item \emph{Pages 1 à 14}, « Structure d'effectivité » : comment le groupe
  de Galois motivique pourrait décrire la sous-catégorie des motifs effectifs,
  par des cocaractères et des cônes ; le cas d'un tore.
  \item \emph{Pages 16 à 20}, « Équivalence d'Albanese et de Picard des cycles
  algébriques » : un seul calcul, dont le signe dépend d'un choix.
  \item \emph{Pages 23 à 38}, « Théorie de l'équivalence d'Albanese des
  cycles », ses paragraphes 1 à 3 : l'anneau des classes à équivalence
  d'Albanese, les familles paramétrées par des variétés abéliennes, les
  homomorphismes $\ell$-adiques qu'elles induisent.
  \item \emph{Pages 39 à 41}, « Astuce de Tate sur les cycles décomposables ».
  \item \emph{Pages 43 à 46}, la lettre à Serre du 27 août 1965 :
  conjectures A, A', A'' et B.
  \item \emph{Pages 47 à 70}, « Sorites sur VA, Picard, Albanese » : un résumé,
  puis ses paragraphes 1 et 2 --- les correspondances agissent sur les
  variétés d'Albanese et de Picard, le théorème de Lefschetz pour $H^{1}$ avec
  un inverse algébrique, sa réciproque ---, et deux questions.
  \item \emph{Pages 72 et 73}, deux feuillets sur Picard, Albanese et la
  cohomologie entière.
  \item \emph{Pages 75 à 88}, une suite qu'il numérote « I » : une
  axiomatique des variétés abéliennes à isogénie près et de leurs cycles.
  \item \emph{Pages 90 et 91}, le tableau des anneaux de classes de cycles.
  \item \emph{Pages 92 à 99}, « Mumford--Nakai / Cycles algébriques amples » :
  le critère numérique d'amplitude, les cônes d'une surface, puis la
  dimension supérieure.
  \item \emph{Pages 101 à 122}, « Conjectures de Lefschetz pour cycles
  algébriques » : les surfaces, les conjectures locales, les conjectures
  globales pour les groupes de Chow, vérifiées en dimension 2, 3 et 4, puis les
  familles limitées de cycles.
  \item \emph{Pages 125 à 127}, une chaîne d'énoncés équivalents sur les cycles
  algébriquement équivalents à zéro.
  \item \emph{Pages 130 et 131}, les groupes de Chow du membre générique d'un
  système linéaire, et quelques notes au crayon.
\end{enumerate}

\subsection*{Les moutures, gardées distinctes}

Deux suites numérotées, celle des pages 23 à 38 et celle des pages 47 à 69,
emploient les mêmes numéros pour des énoncés différents : il y a un Théorème
1.6 à la page 29 et un autre à la page 55, un Corollaire 2.2 à la page 33 et un
autre à la page 61, une Proposition 2.7 à la page 34 et une autre à la page 67
--- sans compter le Corollaire 2.2 barré de la page 60. Toute citation porte
donc ici sa page.

Les conjectures de type Lefschetz viennent trois fois, et la différence est du
contenu : dans la lettre (pages 43 à 46), pour les cycles à équivalence
algébrique près tensorisés par $\mathbb{Q}$ ; sur un anneau local (pages 105 à
108) ; et pour les groupes de Chow, à équivalence rationnelle, algébrique et
$\tau$ près (pages 110 à 115). La lettre énonce des conjectures ; les pages
105 à 115 en dressent des tableaux et les essaient.

\subsection*{Ce que seule une lecture d'ensemble peut dire}

Cinq liens, qu'aucun feuillet n'énonce, et qui sont de nous.

La référence que le Théorème 1.4 de la page 28 laisse en blanc --- « résulte de
[ , 1.2.] » --- est comblée par le Corollaire 1.2 de la page 51 : une
correspondance envoie un $0$-cycle de somme d'Albanese nulle sur un $0$-cycle
de somme d'Albanese nulle.

La Remarque 1.5 de la page 28 (« on ignore, déjà pour les $0$-cycles, si
l'équivalence d'Albanese est identique à l'équivalence linéaire »), la note de
la page 29, l'hypothèse 2) de la page 117, la chaîne de la page 125 et les
réserves « si équiv. … $=$ équiv. d'Albanese » des pages 114 et 115
sont une seule question. Mumford y a répondu en 1969, par la négative sur
$\mathbb{C}$ dès que la surface porte une 2-forme holomorphe. Les « ok » de la
page 114 pour les $0$-cycles ne valent que sous l'hypothèse qu'il a réfutée.

Le cycle $U$ dont la page 31 dit qu'« a priori on ne sait pas si on peut le
trouver » est la Question 1) de la page 70.

Le Théorème 3.1 de la page 37 est exactement ce dont aurait besoin la
vérification laissée ouverte à la page 36 --- que la dualité entre familles de
cycles ne dépende que de l'application et non du cycle qui la définit.

Enfin, les « ? » que les pages 114 et 115 marquent « question numérique » sont
exactement les cas où les exemples de Griffiths (1969) et de Clemens (1983)
mettent les énoncés en défaut ; ce sont les mêmes exemples qui réfutent les
conjectures A et A' de la lettre sous leur forme de 1965.

\subsection*{Les conventions, valables pour tout le dossier}

\begin{itemize}
  \item $k$ est un corps algébriquement clos, $\ell$ un nombre premier distinct
  de sa caractéristique $p$ ; $H^{i}(X)$ est la cohomologie étale $\ell$-adique,
  à coefficients $\mathbb{Q}_{\ell}$ sauf mention de $\mathbb{Z}_{\ell}$, et
  $(r)$ la torsion à la Tate. Sauf mention contraire, $X$ et $Y$ sont
  projectives, lisses et connexes, de dimensions $n$ et $m$ ; aux pages 102 à
  122, $Y$ est une section hyperplane lisse de $X$.
  \item $Z^{i}(X)$ est le groupe des cycles de codimension $i$ et
  $\mathrm{CH}^{i}(X)$ le groupe de Chow, à équivalence rationnelle près ---
  ce qu'il appelle, selon l'usage de 1965, « équivalence linéaire », et qu'il
  note $A^{*}(X)$ puis $\mathcal{A}^{*}_{r}(X)$. Pour les autres équivalences,
  $A^{i}_{\sim}(X) = Z^{i}(X)/\!\sim$ avec $\sim$ égal à $\mathrm{alb}$,
  $\mathrm{alg}$, $\tau$, $\mathrm{hom}$ (équivalence homologique
  $\ell$-adique) ou $\mathrm{num}$. Un indice inférieur est une dimension :
  $\mathrm{CH}_{i}$, $A_{i,\mathrm{alg}}$. Le sous-groupe des classes
  équivalentes à zéro pour une équivalence plus grossière s'écrit en indice
  après la parenthèse : $\mathrm{CH}^{i}(X)_{\mathrm{alg}}$ (ses $J^{i}$,
  $\mathcal{A}^{i}_{r}(X)^{a}$), $\mathrm{CH}^{i}(X)_{\tau}$ (ses
  $\mathcal{A}^{i}_{r}(X)^{\tau}$), $A^{i}_{\mathrm{alb}}(X)_{\mathrm{alg}}$ (ses
  $P^{i}(X)$ des pages 33 à 36). Un cycle est $\tau$-équivalent à zéro si l'un
  de ses multiples non nuls est algébriquement équivalent à zéro.
  \item $\mathrm{Alb}_{X}$ est la variété d'Albanese (son $A_{X}$),
  $\mathrm{alb}_{X}$ la somme d'Albanese des $0$-cycles de degré $0$ (son
  $S_{X}$) ; $\mathrm{Pic}^{0}_{X}$ est la variété de Picard, c'est-à-dire la
  composante neutre réduite du schéma de Picard (son $B_{X}$, qu'il note
  $\underline{\mathrm{Pic}}^{00}_{X/k}$). $\mathrm{Alb}_{X}$ est la duale de
  $\mathrm{Pic}^{0}_{X}$, et l'on note $\hat{A}$ la variété duale de $A$ (son
  $A^{*}$), comme au dossier 12.
  \item $\varphi_{X} = \mathrm{alb}_{X}^{*} \colon H^{1}(\mathrm{Alb}_{X}) \to
  H^{1}(X)$, et $\psi_{X} \colon H^{2n-1}(X)(n-1) \to H^{1}(\mathrm{Pic}^{0}_{X})$
  l'homomorphisme défini par la classe de Poincaré ; ce sont des isomorphismes
  (pages 72 et 73).
  \item Pour un cycle $Z$ sur $X \times Y$, on note $Z_{*}$ son action de $Y$
  vers $X$, $Z_{*}(t) = \mathrm{pr}_{X*}(Z \cdot \mathrm{pr}_{Y}^{*} t)$, sur les
  groupes de Chow comme en cohomologie, et ${}^{t}Z_{*}$ l'action de $X$ vers
  $Y$. La page écrit $Z(t)$ et ${}^{t}Z(D)$ sur les cycles, mais
  $a(\gamma(Z))$ en cohomologie dans les deux sens (de $X$ vers $Y$ à la page
  55, de $Y$ vers $X$ à la page 58) ; on normalise.
  \item Aux pages 1 à 14 : le motif de Tate $\mathbb{Q}(-1)$ est effectif et de
  poids $2$ (la page écrit $\Lambda(-1)$), $\varepsilon \colon G \to
  \mathbb{G}_{m}$ est son caractère, $y$ le cocaractère des poids et $i_{1}$ le
  cocaractère de Hodge, avec $\langle y, \varepsilon\rangle = 2$ et $\langle
  i_{1}, \varepsilon\rangle = 1$ --- les conventions du dossier 12.
  \item Aux pages 95 à 99, les cônes portent leurs noms d'aujourd'hui :
  $\mathrm{Eff}$ (effectif), $\mathrm{Pos}$ (positif), $\mathrm{Amp}$ (ample),
  $\mathrm{Nef}^{i} = \mathrm{Eff}_{i}^{\vee}$ ; il écrit $P^{+}$, $Q^{+}$, $R$
  et $P_{i}^{+\circ}$.
\end{itemize}

\subsection*{Ce que le dossier suppose sans l'établir, et ce qu'il annonce sans
le faire}

Les pages 1 à 14 supposent une catégorie de motifs tannakienne et
semi-simple, ce que les conjectures standard donneraient. Les conjectures des
pages 43 à 46 et 101 à 117 sont des conjectures, et le dossier le dit ; les
résultats démontrés (pages 23 à 36, 47 à 69, 72 et 73, 93 à 95) ne dépendent
d'aucune d'elles.

Il annonce et ne fait pas : les items 1° à 3° de la page 2, qui ne sont pas
dans le dossier ; la « condition de saturation » de la page 8 ; une définition
de l'équivalence de Picard, nommée aux pages 16 et 91 ; la vérification du
déplacement des cycles de la Proposition 1.2 (iii) (« à vérifier », page 27) ;
l'indépendance vis-à-vis du cycle, dans la dualité de la page 36 ; le Théorème
3.1 de la page 37 (« dém.\ pas faite ») et sa réciproque, et la démonstration
du paragraphe 3, qui s'arrête page 38 ; le Lemme 1.5 de la page 54 (« nous nous
dispensons ici d'expliciter ») et la vérification de la Proposition 2.3 de la
page 61 ; les dénominateurs de la Proposition 2.4 ; les questions de signe des
pages 61 et 77 ; les deux questions de la page 70 ; le problème (4) et les
« Pb » des pages 106 à 108 ; la démonstration du Théorème 4 de la page 113 ;
une « partie primitive » de $\mathrm{CH}^{*}(X)$ (page 116) ; l'implication
(ii) $\Rightarrow$ (iii) de la page 119, interrompue.

\pagerange{1}{14}
\section*{I. Structures d'effectivité (pages 1 à 14)}

Le feuillet de couverture porte, de sa main et souligné, « Structure
d'effectivité » ; au-dessus, dans un cadre barré de diagonales, un premier
titre : « Les éléments de structure motiviques et leur interprétation en termes
de groupes algébriques sur $\mathbb{Q}$ ». La suite s'ouvre à l'item 4° : les
items 1° à 3° ne sont pas dans le dossier. Les renvois de la page 10
(« tenant compte de 3° », les polarisations modulo $G(\mathbb{R})$) montrent
que 3° traitait de la structure de Hodge et de la polarisation, décrites par un
cocaractère --- l'objet des pages 3 à 11 du dossier 13 ; rien ne permet de dire
si les items manquants sont ces pages-là.

On se place dans le cadre que ces pages supposent sans le dire : une catégorie
$\mathcal{M}$ de motifs sur $\mathbb{Q}$, tannakienne et semi-simple, munie
d'un foncteur fibre ; pour une sous-catégorie tensorielle de type fini, son
groupe de Galois motivique $G$, groupe algébrique réductif sur $\mathbb{Q}$, de
sorte que $\mathcal{M} \simeq \operatorname{Rep}_{\mathbb{Q}}(G)$ ; le motif
de Tate $\mathbb{Q}(-1)$ et son caractère $\varepsilon \colon G \to
\mathbb{G}_{m}$.\footnote{Rien de ce cadre n'est écrit sur les pages 2 à 14 ;
c'est celui des dossiers 12 et 13, qu'elles supposent connu. La semi-simplicité
et l'existence de la catégorie sont conditionnelles aux conjectures standard.
La page écrit $\Lambda(-1)$ pour le motif de Tate.}

\pagerange{2}{4}
\subsection*{L'axiome 4°, et les cocaractères effectifs (pages 2 à 4)}

Une \emph{structure d'effectivité} sur $\mathcal{M}$ est la donnée d'une
sous-catégorie strictement pleine $\mathcal{M}^{+}$, les « motifs effectifs »,
stable par produit tensoriel, par sommes directes et par sous-objets, qui
contient $\mathbb{Q}(-1)$ et telle que, pour tout motif $E$, il existe un
entier $n$ avec $E(-n) = E \otimes \mathbb{Q}(-1)^{\otimes n}$
effectif.\footnote{Devant « sous-objets », deux mots biffés, illisibles. Dans
une catégorie semi-simple, sous-objets, quotients et facteurs directs sont les
mêmes objets, de sorte que la stabilité par sous-objets est la stabilité par
facteurs directs.}

Comme $\mathcal{M}$ est semi-simple, $\mathcal{M}^{+}$ est déterminée par
l'ensemble des classes d'objets simples qu'elle contient. La donnée équivaut
donc à celle d'un cône $K^{+}$ dans l'anneau $K(G)$ des caractères virtuels de
$G$ sur $\mathbb{Q}$ : un sous-monoïde contenant $0$, stable par
multiplication, et engendré par les éléments de la base canonique --- les
classes des représentations irréductibles --- qu'il contient ; la stabilité
par produit dit que les constituants d'un produit de deux objets effectifs
simples sont effectifs. Une telle structure sur $G$ en définit une sur
$G_{\mathbb{C}}$, en prenant pour objets effectifs les constituants des
complexifiés des objets effectifs ; et une structure sur $G_{\mathbb{C}}$
redescend, pourvu qu'elle soit stable par les conjugaisons de
$\operatorname{Aut}(\mathbb{C})$.\footnote{La page dit « et réciproquement »,
sans la réserve. Une sous-catégorie de $\operatorname{Rep}_{\mathbb{C}}(G)$ qui
provient de $\mathbb{Q}$ est stable par l'action de
$\operatorname{Aut}(\mathbb{C})$ sur les représentations ; c'est aussi ce que
la page 6 retrouvera sous la forme d'un cône stable par $\pi_{1}$. Le
« $G_{\mathbb{Q}}$ » de la page est lu sous réserve.}

Il prévoit d'autres axiomes : les motifs effectifs doivent être à poids
positifs ; la structure doit provenir d'une structure du même type sur la
composante neutre $G^{\circ}$ ; et, avec un « (?!) », il ne doit y avoir, pour
chaque entier $i$, qu'un nombre fini de types de motifs effectifs de poids
$i$.\footnote{L'axiome est encadré, d'une écriture plus rapide, et sa fin n'est
pas lue. On le lit pour les motifs \emph{simples} d'une sous-catégorie de type
fini, décrite par un seul groupe algébrique $G$ : pour la catégorie de tous les
motifs, il est faux, puisque les courbes elliptiques sur $\mathbb{Q}$ deux à
deux non isogènes, en nombre infini, donnent autant de motifs simples effectifs
de poids $1$. La page 14 reviendra sur cette finitude pour un tore.}

Un cocaractère $i \colon \mathbb{G}_{m,\mathbb{C}} \to G_{\mathbb{C}}$ est
\emph{effectif} si, pour tout motif effectif $E$, la graduation de
$E_{\mathbb{C}}$ qu'il définit est à degrés positifs. Ses conjugués par
$G(\mathbb{C})$ le sont alors aussi, et son conjugué complexe $\bar{\imath}$.
Réciproquement, un cocaractère $i$ étant donné, soit $\mathrm{Pl}(i)$ la
sous-catégorie pleine des $E$ tels que la $i$-graduation de $E_{\mathbb{C}}$
soit à degrés positifs. Notons $\langle i, \varepsilon\rangle$ l'entier $\nu$
tel que $\varepsilon(i(\lambda)) = \lambda^{\nu}$. Alors $\mathrm{Pl}(i)$ est
strictement pleine, stable par produit tensoriel, sommes et sous-objets ; elle
contient $\mathbb{Q}(-1)$ si $\langle i, \varepsilon\rangle \geqslant 0$ ; et
elle satisfait la dernière condition de 4° si
\[
\langle i, \varepsilon\rangle > 0 ,
\]
car les degrés de $E(-n)$ sont ceux de $E$ augmentés de $n\langle i,
\varepsilon\rangle$.\footnote{La page écrit ici « $E(n) \in
\mathrm{Ob}\,\mathrm{Pl}(i)$ », là où la page 2 écrivait $E(-n)$ ; c'est $E(-n)$
qui convient, $\varepsilon$ étant le caractère de $\mathbb{Q}(-1)$. Au-dessus
des deux $\langle i, \varepsilon\rangle$, il écrit aussi « $\varepsilon(i)$ ».}
Comme les objets de $\mathcal{M}$ sont définis sur $\mathbb{Q}$, la
sous-catégorie $\mathrm{Pl}(i)$ ne change pas quand on remplace $i$ par un
conjugué sous $G(\mathbb{C})$ ou sous $\operatorname{Aut}(\mathbb{C})$ : un
automorphisme $\sigma$ de $\mathbb{C}$ envoie la partie de degré $d$ pour $i$
sur la partie de degré $d$ pour $\sigma(i)$.\footnote{Cette remarque est de
nous ; elle éclaire la page 10.}

Une intersection d'un nombre fini de catégories $\mathrm{Pl}(i)$, avec des $i$
tels que $\langle i, \varepsilon\rangle > 0$, satisfait donc les axiomes de
4°.\footnote{La page dit « par intersection de tels foncteurs, etc. », sans
limiter le nombre. Pour une famille infinie, la dernière condition de 4°
demande une uniformité en $i$ ; c'est ce que la page 8 exprimera par le mot
« strictement ».}

\pagerange{6}{10}
\subsection*{Les cônes de cocaractères, et deux questions (pages 6 à 10)}

Supposons $G$ connexe.\footnote{Hypothèse ajoutée : pour $G$ non connexe, les
classes de conjugaison de cocaractères de $G$ sont celles de $G^{\circ}$ modulo
l'action de $\pi_{0}(G)$, et la dernière égalité ci-dessous doit être
corrigée d'autant.} Soit $T$ un tore maximal de $G_{\mathbb{C}}$, $\Gamma =
\operatorname{Hom}(\mathbb{G}_{m,\mathbb{C}}, T)$ son groupe de cocaractères,
$W = N(T)/T$ le groupe de Weyl et $W' = \operatorname{Aut}(G, T)/T$ ; alors $W
\subset W'$ et $W'/W = \operatorname{Out}(G)(\mathbb{C})$, et $W'$ opère sur
$\Gamma$. Les cocaractères effectifs contenus dans $T$ forment un cône convexe
$\Gamma^{+}$ : c'est le cône dual de l'ensemble des poids de $T$ dans les
motifs effectifs. Il contient $0$, il est invariant par $W$ --- mais peut-être
pas par $W'$ --- et contenu dans le demi-espace $\langle \cdot,
\varepsilon\rangle \geqslant 0$. Il définit un « cône » dans
\[
\Gamma' = \Gamma/W
= \bigl[\operatorname{Hom}(\mathbb{G}_{m}, T)/(N(T)/T)\bigr](\mathbb{C})
= \bigl[\operatorname{Hom}(\mathbb{G}_{m}, G)/G\bigr](\mathbb{C}) ,
\]
l'ensemble des classes de conjugaison de cocaractères. C'est l'ensemble des
points complexes d'un schéma $S$ sur $\mathbb{Q}$, somme de spectres de corps
de nombres, de sorte que $\pi_{1} = \operatorname{Gal}(\bar{\mathbb{Q}}/
\mathbb{Q})$ opère sur $\Gamma' = S(\mathbb{C}) = S(\bar{\mathbb{Q}})$, à
travers un quotient fini ; et le cône de $\Gamma'$ défini par $\mathcal{M}^{+}$
est stable par $\pi_{1}$.\footnote{La page écrit que $\pi_{1}$ opère sur
$\Gamma = S(\mathbb{C})$, sans prime : c'est $\Gamma'$. Le quotient fini est
celui par lequel $\pi_{1}$ agit sur la donnée radicielle basée de $G$. Il note
$\pi_{1}(\mathbb{Q}, \mathbb{C})$ le groupe fondamental de
$\operatorname{Spec}\mathbb{Q}$ en le point géométrique défini par
$\mathbb{C}$.}

Si l'on se borne aux structures d'effectivité qui sont des intersections de
catégories $\mathrm{Pl}(i)$, la donnée d'une telle structure équivaut à celle
d'un cône de $\Gamma'$ stable par $\pi_{1}$, contenu « strictement » dans le
demi-espace défini par $\varepsilon$ --- en marge : « i.e.\ $\varepsilon$ est
intérieur au cône polaire » ---, et satisfaisant une condition de saturation
qu'il ne précise pas (« condition automatiquement satisfaite ? »). La marge dit
bien ce que « strictement » doit vouloir dire : tout cocaractère non nul du
cône a $\langle i, \varepsilon\rangle > 0$, ce qui, pour un cône polyédral
fermé, revient à ce que $\varepsilon$ soit intérieur au cône dual. Quant à la
saturation, elle est automatique pour les cônes polyédraux rationnels fermés
invariants par $W$.\footnote{Vérification de nous. Soit $C \subset
\Gamma_{\mathbb{R}}$ un tel cône et $C^{\vee}$ son dual, invariant par $W$. Une
représentation irréductible a tous ses poids dans l'enveloppe convexe de
l'orbite sous $W$ de son plus haut poids ; elle est donc dans $\bigcap_{i \in
C}\mathrm{Pl}(i)$ si et seulement si son plus haut poids est dans $C^{\vee}$,
et les poids qui apparaissent sont exactement les points entiers de
$C^{\vee}$. Le cône des cocaractères effectifs pour cette intersection est le
dual de ces points entiers, qui engendrent $C^{\vee}$ : c'est $C$ lui-même.}

Le cas le plus intéressant est celui où $\mathcal{M}^{+} =
\mathrm{Pl}(i_{1})$, avec $i_{1}$ le cocaractère de Hodge de 3° : le cône
invariant engendré par $i_{1}$ est contenu dans le demi-espace de
$\varepsilon$, puisque $\varepsilon \circ i_{1} = \mathrm{id}$, c'est-à-dire
$\langle i_{1}, \varepsilon\rangle = 1$.\footnote{La page écrit que le cône
invariant engendré contient bien « $j$ » ; au dossier 13, $j$ désigne le
morphisme $\mathbb{S} \to G_{\mathbb{R}}$ dont le complexifié donne les
cocaractères $i_{1}$, $i_{2}$.} On obtient ainsi une application
\[
I(G) \longrightarrow \mathrm{Eff}(G)
\]
de l'ensemble des cocaractères du type de 3° vers l'ensemble des structures
d'effectivité, et, en tenant compte de la polarisation,
\[
I^{*}(G) \longrightarrow \mathrm{Pol}(G) \times \mathrm{Eff}(G) ,
\]
où $I^{*}(G)$ est l'ensemble des cocaractères qui définissent une
polarisation, modulo l'action de $G(\mathbb{R})$.

Deux questions. Ces applications sont-elles injectives ? Les structures
d'effectivité qui proviennent de la théorie des motifs sont-elles toujours
décrites par un cocaractère, c'est-à-dire dans l'image de $I(G)$ ? En marge :
cette propriété est intrinsèque à la catégorie $\mathcal{M}$, indépendante des
foncteurs fibres. Il propose de regarder $G$ commutatif, et $G =
\operatorname{Gal}(\bar{\mathbb{Q}}/\mathbb{Q})$ --- les motifs d'Artin,
où tous les cocaractères sont triviaux et tous les motifs effectifs ---, de
tordre par un torseur, et d'examiner le cas d'un motif attaché à un point de
$S(\mathbb{Q})$ au moyen du foncteur de De Rham. Des propriétés arithmétiques
plus profondes, conclut-il, viendraient de l'introduction de filtrations dans
l'interprétation des structures 3° et 4°.\footnote{Sur la première question :
comme $\mathrm{Pl}(i)$ ne change pas quand on remplace $i$ par un conjugué
galoisien (page 4), l'application $I(G) \to \mathrm{Eff}(G)$ ne peut être
injective que sur les orbites de $\pi_{1}$ ; nous ne savons pas si elle l'est
sur ces orbites. Sur la seconde : tout motif effectif a des types de Hodge
$(p, q)$ avec $p, q \geqslant 0$, de sorte que $\mathcal{M}^{+} \subset
\mathrm{Pl}(i_{1})$ ; l'égalité est ce que prédit, jointe à la conjecture de
Hodge, la conjecture de Hodge généralisée sous la forme que Grothendieck lui
donnera en 1969. Elle est ouverte. Un paragraphe barré de la page 12 associait
à un point de $S(k)$ deux foncteurs fibres, de Hodge et de De Rham.}

\pagerange{12}{14}
\subsection*{Le cas d'un tore (pages 12 à 14)}

Soit $T$ un tore sur $\mathbb{Q}$, décrit par son groupe de caractères $M$,
réseau muni d'une action continue de $\pi_{1}$, qui se fait à travers un
quotient fini $\pi$. On a $\varepsilon \in M^{\pi}$, le cocaractère des poids
$y \in \check{M}^{\pi}$ et $\langle y, \varepsilon\rangle = 2$.\footnote{La page
note ce cocaractère $i$ ; invariant par $\pi$ et de degré $2$ sur
$\varepsilon$, c'est le cocaractère des poids, que le dossier 12 note $y$.} Les
représentations de $T$ correspondent aux éléments invariants par $\pi$ de
$\mathbb{Z}[M]^{+}$, les combinaisons à coefficients positifs de caractères.

La structure de polarisation ne dépend que du tore réel $T_{\mathbb{R}}$. Il est
extension d'un tore déployé $S$ par un tore compact $U$, c'est-à-dire tel que
$U(\mathbb{R})$ soit compact, et $\varepsilon$ est trivial sur $U$, qui n'a pas
de caractère réel non trivial ; on est ramené à $S$. Grothendieck affirme
qu'on doit avoir $\dim S = 1$, autrement dit que $\operatorname{Ker}
\varepsilon$ est compact. C'est vrai dès que la structure est polarisée, et la
démonstration tient en une ligne.\footnote{La tentative de la page, qui cherche
une contradiction dans un tore déployé de dimension $2$, est barrée de deux
grands traits. Voici l'argument (de nous). Soit $\psi$ une polarisation d'une
représentation fidèle $V$ ; $T$ agit sur $\psi$ par une puissance de
$\varepsilon$, et l'opérateur de Weil $C = h(\sqrt{-1})$ est dans
$T(\mathbb{R})$. Les éléments de $(\operatorname{Ker}\varepsilon)(\mathbb{R})$
commutent à $C$ et préservent $\psi$, donc la forme symétrique définie positive
$\psi(x, Cy)$ ; ils forment un sous-groupe fermé d'un groupe orthogonal
compact. Un tore réel de points réels compacts est anisotrope. C'est, pour un
tore, la condition qu'on impose aujourd'hui aux groupes de Mumford--Tate des
structures de Hodge polarisables ; les tores ainsi munis d'un poids et d'un
cocaractère de Hodge sont de la famille de ceux que classe le groupe de Serre
(Serre, 1968), dont Deligne (1982) a montré qu'il est le groupe de Galois
motivique des motifs de type CM sur $\bar{\mathbb{Q}}$, pour les cycles de Hodge
absolus.}

Page 14, il détermine l'anneau des caractères. Sur $\mathbb{C}$, $T_{\mathbb{C}}
= D_{\mathbb{C}}(M)$ ; l'anneau $\mathbb{Z}[M]$ des caractères virtuels de
$T_{\mathbb{C}}$ porte l'action de $\pi$, et $K(T) = \mathbb{Z}[M]^{\pi}$ est
une $\mathbb{Z}$-algèbre de type fini. Mieux : $\mathbb{Z}[M]$ est l'anneau
affine du tore déployé $\mathbf{D} = D_{\mathbb{Z}}(M)$, sur lequel $\pi$ agit
par automorphismes de groupe, et $K(T)$ celui du quotient $\mathbf{D}/\pi$.
Les caractères effectifs seraient-ils ceux qui prennent des valeurs positives
sur la composante neutre de $\mathbf{D}(\mathbb{R})$ ? Il le demande (« ?? »)
et demande comment définir une structure d'effectivité intrinsèque.\footnote{Le
titre de ce passage, « Détermination du groupe des relations rationnelles sur
$\mathbb{Q}$ », est lu sous réserve pour ses deux mots centraux.}

Le procédé standard est de se fixer un cocaractère $i_{1} \colon
\mathbb{G}_{m,\mathbb{C}} \to T_{\mathbb{C}}$, c'est-à-dire $i_{1}^{*} \colon M
\to \mathbb{Z}$, d'où $\mathbb{Z}[M] \to \mathbb{Z}[t, t^{-1}]$ et un morphisme
$\alpha \colon \mathbb{G}_{m,\mathbb{Z}} \to \mathbf{D}$, non compatible à
$\pi$. Un caractère $\chi \in K(T)$ est effectif si $\chi \circ \alpha$ est un
polynôme en $t$, c'est-à-dire reste fini sur $\mathbb{R}^{*}$ quand $t \to 0$.
Il en tire, dans un dernier paragraphe très serré, que si $\dim T \geqslant 3$
la finitude des caractères effectifs de poids donné est
fausse.\footnote{La lecture de ce paragraphe est très incertaine, et le
raisonnement n'est pas sur la page. Voici ce qu'on peut vérifier (de nous). Un
caractère $\chi$ invariant par $\pi$ a un support stable par $\pi$, de sorte que
la condition porte automatiquement sur tous les conjugués $\sigma(i_{1})$ ; si,
comme pour un cocaractère de Hodge, le conjugué complexe de $i_{1}$ est $y -
i_{1}$, les caractères effectifs de poids $w$ sont les points entiers du
polyèdre défini par $0 \leqslant \langle m, \sigma(i_{1})\rangle \leqslant w$.
Ce polyèdre est borné, et l'ensemble fini, si et seulement si les conjugués de
$i_{1}$ engendrent $\check{M} \otimes \mathbb{Q}$, c'est-à-dire si $T$ est le
plus petit sous-tore défini sur $\mathbb{Q}$ par lequel passe $i_{1}$ --- son
tore de Mumford--Tate. Dans ce cas la finitude vaut en toute dimension ;
sinon elle tombe, déjà en dimension $2$. Si la page vise la positivité sur
$\mathbf{D}(\mathbb{R})^{\circ}$ du paragraphe précédent, nous n'avons pas
reconstitué son argument.}

\pagerange{16}{20}
\section*{II. Équivalence d'Albanese et de Picard : un exemple (pages 16 à 20)}

Le feuillet 16 porte le titre « Équivalence d'Albanese et de Picard des cycles
algébriques ». L'équivalence d'Albanese sera définie aux pages 23 et suivantes ;
l'équivalence de Picard n'est définie nulle part dans le dossier. Les pages 18
et 20, sous une grande lettre « E », font un seul calcul.

Soit $A$ une courbe elliptique, $P$ une courbe lisse, complète et connexe, et $Y
= A \times P$. Sur $A \times A$, le diviseur $\Theta_{0} = \Delta - A \times
\{0\} - \{0\} \times A$ définit l'autodualité de $A$ ; sur $A \times Y = A
\times A \times P$, on pose $\Theta = \Theta_{0} \times P$. Soit $L$ un faisceau
inversible sur $Y$, de classe $c$, et $X$ la complétion projective du fibré en
droites le long de la section nulle duquel le fibré normal est $L$ : une
variété de dimension $3$, avec sa section nulle $i \colon Y \to X$. Le cycle
\[
z = (\mathrm{id}_{A} \times i)_{*}(\Theta)
\]
est de codimension $2$ sur $A \times X$ : c'est une famille, paramétrée par
$A$, de courbes $z_{a} = i_{*}\bigl((\{a\} - \{0\}) \times P\bigr)$ sur $X$,
algébriquement équivalentes à zéro. Soit $s \colon A \times X \to A$ la
projection ; on calcule le degré du $0$-cycle $s_{*}(z^{2})$.

Posons $i' = \mathrm{id}_{A} \times i$. La formule d'auto-intersection donne
$z^{2} = i'_{*}(\Theta^{2} \cdot c')$, où $c' = i'^{*}i'_{*}(1) = 1_{A} \times
c$ est la classe du fibré normal. D'où $s_{*}(z^{2}) = t_{*}(\Theta^{2}\cdot
c')$, avec $t \colon A \times Y \to A$ la projection, et, par la formule de
projection pour $p \colon A \times Y \to Y$,
\[
\deg s_{*}(z^{2}) = \deg\bigl(p_{*}(\Theta^{2}) \cdot c\bigr) .
\]
Comme $p = p' \times \mathrm{id}_{P}$, avec $p' \colon A \times A \to A$ la
seconde projection, $p_{*}(\Theta^{2}) = p'_{*}(\Theta_{0}^{2}) \times P$. Or
\[
\Theta_{0}^{2} = \Delta^{2} - 2(0,0) - 2(0,0) + 2(0,0) = -2\,(0,0) ,
\]
car $\Delta^{2} = 0$ : le fibré tangent de $A$ est trivial. Donc
$p_{*}(\Theta^{2}) = -2\,[\{0\} \times P]$ et
\[
\deg s_{*}(z^{2}) = -2\,\deg\bigl(c \cdot [\{0\} \times P]\bigr) .
\]
Pour $c$ la classe de $A \times \{a\}$, $a \in P(k)$, on trouve $-2$ ;
remplaçant $c$ par $-c$, on trouve $+2$ : « des signes différents
!! ».\footnote{La lettre notée ici $\deg$ est lue « $K$ » sous réserve ; le
calcul n'a de sens que si c'est le degré. La page écrit $p_{*}(\Theta^{2}) =
-2(0 \times 0 \times P)$ et $c = A \times A \times (a)$, là où les cycles vivent
sur $Y = A \times P$ ; elle justifie $\Delta^{2} = 0$ « à équivalence
algébrique près, puisque $\chi(A) = 0$ » : c'est vrai à équivalence
rationnelle, par la trivialité du fibré tangent. Le diagramme de la marge
dessine la flèche $i'$ à rebours.}

Ce que le calcul met en évidence est, nous semble-t-il, ceci. Le fibré
projectif $\mathbb{P}(\mathcal{O}_{Y} \oplus L)$ est isomorphe à
$\mathbb{P}(\mathcal{O}_{Y} \oplus L^{-1})$ ; il a deux sections, de fibrés
normaux $L$ et $L^{-1}$. Les deux signes sont donc obtenus sur \emph{une même}
variété de dimension $3$, par deux familles de courbes algébriquement
équivalentes à zéro. La quantité $z \mapsto \deg s_{*}(z^{2})$ n'a pas de signe
fixe : aucune classe ample n'intervient, et, sur le $H^{3}$ d'une variété de
dimension $3$, la forme d'intersection n'est définie qu'en restriction à la
partie primitive relativement à une classe ample. C'est ce que la conjecture B
de la lettre des pages 43 à 46 incorporera.\footnote{Cette lecture est de nous ;
la page s'arrête sur le constat.}

\pagerange{23}{38}
\section*{III. Théorie de l'équivalence d'Albanese (pages 23 à 38)}

Une rédaction au propre, sous le titre souligné « Théorie de l'équivalence
d'Albanese des cycles », en paragraphes numérotés 1 à 3, non paginée ; en
marge de la première page : « Consulter Weil, Critères
d'équivalence ».\footnote{Sans doute A. Weil, « Sur les critères d'équivalence
en géométrie algébrique », \emph{Math.\ Annalen} 128 (1954).} Elle s'arrête aux
deux tiers de la page 38, au milieu d'une démonstration.

\pagerange{23}{29}
\subsection*{1. L'anneau des classes à équivalence d'Albanese (pages 23 à 29)}

\textbf{Définition 1.1.} Soit $X$ un schéma lisse sur $k$ et $\mathfrak{z}$ un
cycle sur $X$. On dit que $\mathfrak{z}$ est \emph{équivalent à zéro au sens
d'Albanese} s'il existe un schéma projectif, lisse et connexe $T$, un
$0$-cycle $\mathfrak{t}$ sur $T$ de degré $0$ et de somme d'Albanese
$\mathrm{alb}_{T}(\mathfrak{t}) = 0$, et un cycle $Z$ sur $X \times T$ tel que
$Z_{*}(\mathfrak{t})$ soit défini et égal à $\mathfrak{z}$.\footnote{« Projectif »
et « connexe » sont ajoutés par lui entre crochets ; le premier est nécessaire à
la somme d'Albanese telle qu'elle sert au Théorème 1.4, comme il le note en
marge de la page 28 (« ici il semble qu'on utilise le fait que dans 1.1 on prend
$T$ projectif »). Il avait d'abord écrit « $X$ projectif », puis biffé
« projectif ».}

\textbf{Proposition 1.2.} (i) Les cycles équivalents à zéro au sens d'Albanese
forment un sous-groupe gradué $N(X)$ de $Z^{*}(X)$, qui contient les cycles
rationnellement équivalents à zéro ; son image $N'(X)$ dans $\mathrm{CH}^{*}(X)$
le détermine, et $Z^{*}(X)/N(X) \simeq \mathrm{CH}^{*}(X)/N'(X)$.
(ii) $N'(X)$ est un idéal, de sorte que $A^{*}_{\mathrm{alb}}(X) =
\mathrm{CH}^{*}(X)/N'(X)$ est un anneau.
(iii) Pour tout morphisme $f \colon X \to Y$ de variétés quasi-projectives
lisses, $f^{*}$ envoie $N'(Y)$ dans $N'(X)$ et induit un homomorphisme d'anneaux
$A^{*}_{\mathrm{alb}}(Y) \to A^{*}_{\mathrm{alb}}(X)$.
(iv) Si de plus $f$ est propre, l'image directe $f_{*}$ envoie $N'(X)$ dans
$N'(Y)$ et induit $A^{*}_{\mathrm{alb}}(X) \to A^{*}_{\mathrm{alb}}(Y)$.
(v) L'équivalence d'Albanese entraîne l'équivalence algébrique ; on peut donc
parler des classes de $A^{*}_{\mathrm{alb}}(X)$ algébriquement équivalentes à
zéro.

Les démonstrations sont formelles. Pour la somme, si $\mathfrak{z} =
Z_{*}(\mathfrak{t})$ et $\mathfrak{z}' = Z'_{*}(\mathfrak{t}')$ avec des
paramètres $T$, $T'$, on remonte $Z$ et $Z'$ en $Z_{1}$, $Z'_{1}$ sur $X \times
T \times T'$, on choisit des points $a \in T$, $a' \in T'$ en lesquels $Z$ et
$Z'$ se spécialisent bien, et l'on pose $\mathfrak{t}_{2} = \mathfrak{t} \times
a' + a \times \mathfrak{t}'$, $Z_{2} = Z_{1} + Z'_{1}$ : les $0$-cycles
$\mathfrak{t} \times a'$ et $a \times \mathfrak{t}'$ sont de somme d'Albanese
nulle sur $T \times T'$, les termes croisés sont nuls, et $Z_{2*}(\mathfrak{t}_{2})
= \mathfrak{z} + \mathfrak{z}'$ ; le passage à $-\mathfrak{z}$ est
immédiat.\footnote{La page annonce « $\mathfrak{z} - \mathfrak{z}'$ » et
calcule $\mathfrak{z} + \mathfrak{z}'$, la ligne sur $-\mathfrak{z}$ ayant été
biffée plus haut. Les points $a$, $a'$ ne sont pas introduits sur la page ; en
marge, « Lemme à dégager ».} L'équivalence rationnelle est contenue parce que
sur $T = \mathbb{P}^{1}$ tout $0$-cycle de degré $0$ a une somme d'Albanese
nulle. Pour (ii), (iii) et (iv), on a respectivement $Z_{*}(\mathfrak{t}) \cdot
w = (Z \cdot \mathrm{pr}_{X}^{*}w)_{*}(\mathfrak{t})$, $f^{*}Z_{*}(\mathfrak{t})
= ((f \times \mathrm{id}_{T})^{*}Z)_{*}(\mathfrak{t})$ et
$f_{*}Z_{*}(\mathfrak{t}) = ((f \times \mathrm{id}_{T})_{*}Z)_{*}(\mathfrak{t})$,
et (v) vient de ce qu'un $0$-cycle de degré $0$ sur $T$ connexe est
algébriquement équivalent à zéro.\footnote{La page raisonne sur les cycles
eux-mêmes, et doit donc déplacer $Z$ et $\mathfrak{t}$ pour que les
intersections soient définies : pour (iii), elle laisse la vérification de
l'existence de bons représentants (« à vérifier ») ; (ii) se ramène par
produits cartésiens à (iii) ; (iv) se fait « par l'AQT » --- abréviation lue
sous réserve, que nous ne reconstituons pas. Dans la théorie de l'intersection
de Fulton (1984), qui opère sur les classes, les trois formules ci-dessus sont
des identités, et aucun déplacement n'est nécessaire.}

\textbf{Scholie 1.3.} Les anneaux $A^{*}_{\mathrm{alb}}(X)$ ont les propriétés
fonctorielles familières des anneaux de Chow ; en particulier la formule de
projection vaut. Là où équivalence rationnelle et équivalence algébrique
coïncident --- le point, les espaces projectifs, les variétés de drapeaux ---,
elles coïncident avec l'équivalence d'Albanese, qui est entre les deux. C'est
une relation d'équivalence adéquate au sens de Samuel.\footnote{P. Samuel,
« Relations d'équivalence en géométrie algébrique », congrès international
d'Édimbourg, 1958. Le mot n'est pas sur la page.}

\textbf{Théorème 1.4.} Soit $X$ projective, lisse et connexe, de dimension $n$.
(i) Un $0$-cycle $\mathfrak{t}$ sur $X$ est équivalent à zéro au sens d'Albanese
si et seulement s'il est de degré $0$ et $\mathrm{alb}_{X}(\mathfrak{t}) = 0$.
(ii) Un diviseur $D$ sur $X$ est équivalent à zéro au sens d'Albanese si et
seulement s'il est rationnellement équivalent à zéro.

Autrement dit, $A_{0,\mathrm{alb}}(X)$ est extension de $\mathbb{Z}$ par
$\mathrm{Alb}_{X}(k)$, et $A^{1}_{\mathrm{alb}}(X) = \mathrm{Pic}(X)$. Pour (i),
un $0$-cycle de degré $0$ et de somme nulle est équivalent à zéro par
définition, avec $T = X$ et $Z = \Delta_{X}$. La réciproque demande qu'une
correspondance $Z$ sur $X \times T$ envoie un $0$-cycle de somme nulle sur un
$0$-cycle de somme nulle : c'est le Corollaire 1.2 de la page 51, qui la déduit
de ce que $Z$ induit un homomorphisme $\mathrm{Alb}_{T} \to \mathrm{Alb}_{X}$
compatible aux sommes.\footnote{La page renvoie à « [ , 1.2.] » et laisse le
premier terme de la référence en blanc. Le rapprochement avec la page 51 est de
nous.} Pour (ii), si $Z$ est un diviseur sur $X \times T$, on a
$Z_{*}(\mathfrak{t}) = u(\mathrm{alb}_{T}(\mathfrak{t}))$ pour l'homomorphisme
$u \colon \mathrm{Alb}_{T} \to \mathrm{Pic}^{0}_{X}$ défini par la
correspondance divisorielle $Z$, qui est nul si $\mathrm{alb}_{T}(\mathfrak{t})
= 0$.

\textbf{Remarque 1.5.} On ignore, déjà pour les $0$-cycles, si l'équivalence
d'Albanese est l'équivalence rationnelle --- et il se demande, avec trois points
d'interrogation, si ce serait vrai pour tous les cycles.\footnote{La réponse est
non. D. Mumford, « Rational equivalence of 0-cycles on surfaces »,
\emph{J. Math.\ Kyoto Univ.} 9 (1969) : sur une surface projective lisse
complexe dont le genre géométrique $p_{g}$ est non nul, le noyau de
$\mathrm{alb}_{X}$ sur les $0$-cycles de degré $0$ est « de dimension infinie »,
en particulier non nul. Pour $p_{g} = 0$, la conjecture de Bloch prédit que ce
noyau est nul ; elle est démontrée pour les surfaces qui ne sont pas de type
général (Bloch, Kas et Lieberman, 1976) et pour certaines de type général, et
ouverte en général. Cette question commande les réserves des pages 114 et 115
et la chaîne de la page 125.}

\textbf{Théorème 1.6 (page 29).} Si deux classes $\mathfrak{z}$,
$\mathfrak{z}'$ de $A^{*}_{\mathrm{alb}}(X)$ sont algébriquement équivalentes à
zéro, leur produit est nul. En effet, l'équivalence algébrique se témoigne par
des courbes : $\mathfrak{z} = Z_{*}((x) - (y))$ et $\mathfrak{z}' =
Z'_{*}((x') - (y'))$ avec des paramètres $T$, $T'$ ; alors $\mathfrak{z}\cdot
\mathfrak{z}'$ est l'image par la correspondance $(\Delta_{X} \times
\mathrm{id})^{*}(Z \times Z')$ du $0$-cycle
\[
\mathfrak{t} = \bigl((x) - (y)\bigr) \times \bigl((x') - (y')\bigr)
= (x, x') - (x, y') - (y, x') + (y, y')
\]
de $T \times T'$, dont la somme d'Albanese est nulle, puisque $\mathrm{Alb}_{T
\times T'} = \mathrm{Alb}_{T} \times \mathrm{Alb}_{T'}$. Il demande entre
crochets s'il suffirait, pour $X$ projective, que l'une des classes soit
algébriquement et l'autre numériquement équivalente à zéro, et note : « on
ignore si cet $\mathfrak{t}$ est linéairement équivalent à zéro, même si $T,
T'$ sont des courbes ! ».\footnote{Non en général. Si $T$ et $T'$ sont des
courbes, le noyau d'Albanese de $T \times T'$ est engendré par ces cycles : pour
$\sum n_{i}(x_{i}, y_{i})$ de degré $0$ et de somme nulle, on écrit $\sum
n_{i}(x_{i}, y_{i}) = \sum n_{i}\,((x_{i}) - (a)) \times ((y_{i}) - (b))$ plus
des termes $(\sum n_{i}x_{i}) \times b$ et $a \times (\sum n_{i}y_{i})$, qui
sont nuls par le théorème d'Abel sur chaque courbe. Or ce noyau est non nul, et
même de dimension infinie, dès que les deux courbes sont de genre $\geqslant 1$
sur $\mathbb{C}$, puisque $p_{g}(T \times T') = g(T)\,g(T')$ (Mumford, 1969).}
Il y aurait intérêt, conclut-il, à faire de ce lemme un énoncé sur le produit
extérieur $\times$ --- ce que fera la page 125.

\pagerange{30}{36}
\subsection*{2. Familles de classes à équivalence d'Albanese près (pages 30 à
36)}

\textbf{Théorème 2.1.} Soit $Z \in A^{i}_{\mathrm{alb}}(X \times Y)$.
L'application $\mathfrak{t} \mapsto Z_{*}(\mathfrak{t})$, de
$A_{0,\mathrm{alb}}(Y)$ dans $A^{i}_{\mathrm{alb}}(X)$, est un homomorphisme ;
comme $\mathrm{alb}_{Y}$ identifie $A_{0,\mathrm{alb}}(Y)$ au groupe des points
de toutes les composantes $\mathrm{Alb}^{d}_{Y}$ ($d \in \mathbb{Z}$), elle
s'écrit $Z_{*}(\mathfrak{t}) = u^{\bullet}_{Z}(\mathrm{alb}_{Y}(\mathfrak{t}))$,
et l'on note $u_{Z} \colon \mathrm{Alb}_{Y}(k) \to A^{i}_{\mathrm{alb}}(X)$ sa
restriction au degré $0$. Alors : (i) $u_{Z}$ ne dépend que de la classe de $Z$
à $\tau$-équivalence près ; (ii) il existe une isogénie $\varphi \colon A' \to
\mathrm{Alb}_{Y}$ et un cycle $Z' \in A^{i}_{\mathrm{alb}}(X \times A')$ tel que
$Z'_{*}((0)) = 0$ --- ce qui revient à dire que $v_{Z'} \colon a' \mapsto
Z'_{*}((a'))$ est un homomorphisme $A'(k) \to A^{i}_{\mathrm{alb}}(X)$ --- et
tel que $u_{Z} \circ \varphi = v_{Z'}$. On peut prendre pour $\varphi$ la
multiplication par un entier sur $\mathrm{Alb}_{Y}$.\footnote{Il note
$\mathrm{Alb}^{*}_{Y}$ le « torseur de tous les degrés » ; en marge de (i) :
« à réserver pour plus tard ».}

Pour (ii), s'il existait un cycle $U$ sur $Y \times \mathrm{Alb}_{Y}$, de
codimension $m$, tel que $\mathrm{alb}_{Y}(U_{*}(a)) = a$ pour tout $a$, le
cycle $Z' = Z \circ U$, normalisé en $0$, conviendrait avec $\varphi =
\mathrm{id}$. On ne sait pas trouver un tel $U$, mais on en trouve un avec
$\mathrm{alb}_{Y}(U_{*}(a)) = na$ pour un entier $n > 0$.\footnote{La page
l'affirme sans justification. C'est classique : une courbe intersection
complète $C \subset Y$ dont la jacobienne se surjecte sur $\mathrm{Alb}_{Y}$,
la réductibilité de Poincaré, et l'inversion de Jacobi sur $C$. L'existence d'un
$U$ avec $n = 1$ est la Question 1) de la page 70.} Posons $Z'(a) =
Z_{*}(U_{*}(a) - U_{*}(0))$. Si $\mathfrak{t}$ est de degré $0$ et $a =
\mathrm{alb}_{Y}(\mathfrak{t})$, les $0$-cycles $U_{*}(a) - U_{*}(0)$ et
$n\mathfrak{t}$ ont même degré et même somme, donc sont équivalents au sens
d'Albanese, et $Z'(a) = n Z_{*}(\mathfrak{t})$ dans
$A^{i}_{\mathrm{alb}}(X)$ ; le même argument montre que $a \mapsto Z'(a)$ est
additif.\footnote{La page écrit « $U(a) = n\mathfrak{t}$ » ; l'égalité n'a lieu
qu'à équivalence d'Albanese près, ce qui suffit puisque $Z_{*}$ passe au
quotient. Elle invoque aussi la divisibilité de $\mathrm{Alb}_{Y}(k)$ pour
écrire tout élément de degré $0$ comme un multiple ; un passage de cinq lignes
est biffé.} D'où $u_{Z} \circ n = v_{Z'}$. Pour (i) : si $Z$ est
algébriquement équivalent à zéro, $u_{Z} = 0$ par le Théorème 1.6 de la page
29, appliqué à $Z \cdot \mathrm{pr}_{Y}^{*}\mathfrak{t}$ ; si $mZ$ l'est, $m
u_{Z} = 0$, donc $u_{Z} = 0$ puisque $\mathrm{Alb}_{Y}(k)$ est divisible.

\textbf{Corollaire 2.2 (page 33).} Un cycle de codimension $i$ sur $X$ est
algébriquement équivalent à zéro si et seulement s'il est équivalent au sens
d'Albanese à un cycle $W_{*}((a_{0}))$, où $A$ est une variété abélienne et $W$
un cycle de codimension $i$ sur $X \times A$ tel que $a \mapsto W_{*}((a))$ soit
un homomorphisme $A(k) \to A^{i}_{\mathrm{alb}}(X)$.\footnote{La page écrit
« Albanese-équivalent à un diviseur de la forme $Z(a_{0})$ » et place d'abord
$Z$ « sur une $A$ », puis, en se corrigeant, sur « $A \times X$ » : c'est un
cycle de codimension $i$, pas un diviseur. Le numéro du corollaire est
surchargé.}

\textbf{Corollaire 2.3.} Le groupe $A^{i}_{\mathrm{alb}}(X)_{\mathrm{alg}}$ des
classes algébriquement équivalentes à zéro est divisible : c'est une réunion
d'images de groupes $A(k)$.

\textbf{Corollaire 2.4.} Si $\mathfrak{z}$ est algébriquement et
$\mathfrak{z}'$ $\tau$-équivalent à zéro, $\mathfrak{z}\mathfrak{z}' = 0$ dans
$A^{*}_{\mathrm{alb}}(X)$.\footnote{La page ne donne pas de démonstration, et la
divisibilité seule ne suffit pas : un élément divisible tué par $m$ n'est pas
nul pour autant. Mais par 2.2, $\mathfrak{z} = W_{*}((a_{0})) = m\,W_{*}((b))$
avec $a_{0} = mb$, et $\mathfrak{z}\mathfrak{z}' = W_{*}((b)) \cdot
(m\mathfrak{z}') = 0$ par le Théorème 1.6 de la page 29.}

\textbf{Définition 2.5.} Une application $\phi \colon Y(k) \to
A^{i}_{\mathrm{alb}}(X)$ est \emph{algébrique} si 1°) son prolongement par
linéarité aux $0$-cycles se factorise par $A_{0,\mathrm{alb}}(Y)$, et 2°)
l'homomorphisme induit sur $\mathrm{Alb}_{Y}(k)$ est tel que, pour un entier $N >
0$, $a \mapsto \phi(Na)$ soit induit par un cycle de $A^{i}_{\mathrm{alb}}(X
\times \mathrm{Alb}_{Y})$. \textbf{Proposition 2.6.} Une application définie
par un cycle de $X \times Y$ est algébrique (c'est 2.1).
\textbf{Proposition 2.7 (page 34).} Les applications algébriques forment un
groupe, et $\phi$ est algébrique si $N\phi$ l'est ; elles sont fonctorielles en
$Y$, et en $X$ par les correspondances à coefficients entiers.

Un point $y_{0}$ étant fixé, les familles algébriques $Y(k) \to
A^{i}_{\mathrm{alb}}(X)$ se décomposent en une constante et un homomorphisme
algébrique $\mathrm{Alb}_{Y}(k) \to A^{i}_{\mathrm{alb}}(X)$. On est donc ramené
aux homomorphismes algébriques $A(k) \to A^{i}_{\mathrm{alb}}(X)$, et, comme
$A(k) = A_{0,\mathrm{alb}}(A)_{\mathrm{alg}}$, à la notion générale :
un homomorphisme $\phi \colon A^{j}_{\mathrm{alb}}(X')_{\mathrm{alg}} \to
A^{i}_{\mathrm{alb}}(X)_{\mathrm{alg}}$ est \emph{algébrique} si $N\phi$ est
défini par un cycle sur $X \times X'$ pour un $N > 0$. Il l'est alors par un
cycle à coefficients rationnels --- mais un cycle à coefficients rationnels
quelconque ne définit pas un tel homomorphisme, « à moins que les degrés soient
les bons ». Ces homomorphismes se composent ; peuvent-ils s'inverser, quand ce
sont des isomorphismes ? Il le demande.

Vient une \emph{dualité}, pour $i + j = n + 1$. Si $A$ et $\hat{A}$ sont deux
variétés abéliennes duales, un homomorphisme $A(k) \to
A^{i}_{\mathrm{alb}}(X)_{\mathrm{alg}}$ défini par $Z \in
A^{i}_{\mathrm{alb}}(X \times A)$ donne, par le transposé,
\[
{}^{t}Z_{*} \colon A^{j}_{\mathrm{alb}}(X)_{\mathrm{alg}} \longrightarrow
A^{i+j-n}_{\mathrm{alb}}(A)_{\mathrm{alg}} = A^{1}_{\mathrm{alb}}(A)_{\mathrm{alg}}
= \mathrm{Pic}^{0}(A) = \hat{A}(k) ,
\]
d'où une correspondance entre $\mathrm{Hom}_{\mathrm{alg}}(A,
A^{i}_{\mathrm{alb}}(X)_{\mathrm{alg}})$ et
$\mathrm{Hom}_{\mathrm{alg}}(A^{j}_{\mathrm{alb}}(X)_{\mathrm{alg}}, \hat{A})$
--- à condition de vérifier que ${}^{t}Z_{*}$ ne dépend que de l'application et
non de $Z$, ce que la page signale et ne fait pas. Il note la question de la
transposition en général, « cf.\ plus bas ».\footnote{L'égalité
$A^{1}_{\mathrm{alb}}(A)_{\mathrm{alg}} = \mathrm{Pic}^{0}(A)$ est le Théorème
1.4 (ii). La vérification manquante est, en cohomologie, exactement ce
qu'énonce le Théorème 3.1 de la page 37 : la codimension $j = n + 1 - i$ est
celle des cycles dont l'invariant continu vit dans $H^{2j-1} = H^{2(n-i)+1}$.}

\pagerange{37}{38}
\subsection*{3. Homomorphismes $\ell$-adiques associés à une famille abélienne
de cycles (pages 37 et 38)}

Soit $A$ une variété abélienne et $Z \in A^{i}_{\mathrm{alb}}(X \times A)$ ; on
a $u_{Z} \colon A(k) \to A^{i}_{\mathrm{alb}}(X)$, mais aussi
\[
{}^{t}Z_{*} \colon H^{2(n-i)+1}(X)(n-i) \longrightarrow H^{1}(A) ,
\]
à coefficients $\mathbb{Z}_{\ell}$. Il affirme qu'à torsion près cet
homomorphisme ne dépend que de $u_{Z}$, et plus précisément :

\textbf{Théorème 3.1.} Si $u_{Z} = 0$, alors ${}^{t}Z_{*} \colon
H^{2(n-i)+1}(X)(n-i) \to H^{1}(A)$ est nul (à coefficients $\mathbb{Q}_{\ell}$).

En marge : « dém.\ pas faite », entouré, et « A-t-on une réciproque ? ». La
démonstration commencée choisit une courbe « ample » $C \to A$, de sorte que
$H^{1}(A) \to H^{1}(C)$ soit injectif et que $C(k)$ engendre $A(k)$ : les deux
nullités se lisent alors sur $C$, et l'on est ramené à $A = C$. La décomposition
de Künneth de $H^{2i}(X \times C)$ est algébrique, puisque les projecteurs de
Künneth d'une courbe le sont ; un cycle qui « appartient » à l'un des deux
facteurs extrêmes, $H^{2i}(X) \otimes H^{0}(C)$ ou $H^{2i-2}(X) \otimes
H^{2}(C)$, donne un $u$ nul et un homomorphisme cohomologique nul. On est
ramené aux cycles du facteur médian, dont la projection sur $X$ et la
restriction à $X$ sont nulles --- et la page s'arrête là.\footnote{Sur
$\mathbb{C}$, l'énoncé se démontre aujourd'hui par l'application
d'Abel--Jacobi de Griffiths (1968) : elle passe au quotient par l'équivalence
d'Albanese, puisqu'elle est compatible aux correspondances et coïncide avec la
somme d'Albanese sur les $0$-cycles ; si $u_{Z} = 0$, l'homomorphisme de tores
complexes $A \to J^{i}(X)$ défini par $Z$ est nul, donc aussi son effet sur
$H_{1}$, et, par dualité de Poincaré, ${}^{t}Z_{*}$. Nous ne savons pas dire si
la forme $\ell$-adique est établie en caractéristique $p$, ni rien de la
réciproque qu'il demande.}

\pagerange{39}{41}
\section*{IV. L'astuce de Tate sur les cycles décomposables (pages 39 à 41)}

Le feuillet 39 porte seul le titre ; le feuillet 40, d'une encre plus noire sur
un papier jauni, et le haut du feuillet 41, qui achève sa dernière phrase,
donnent l'argument. Un cycle est ici \emph{décomposable} s'il est produit de
classes de diviseurs ; on note $D^{2i}(A, \mathbb{Q}_{\ell}(i))$ le sous-espace
de $H^{2i}(A, \mathbb{Q}_{\ell}(i))$ engendré par les produits de $i$ classes de
diviseurs.

Soit $A$ une variété abélienne ; son module de Tate $T_{\ell}(A)$ est un module
sur ce qu'il appelle la « pseudo-algèbre de Lie » $\mathfrak{g}$, et l'on pose
$T_{K}(A) = T_{\ell}(A) \otimes K$ pour une clôture algébrique $K$ de
$\mathbb{Q}_{\ell}$, module sur $\mathfrak{g}_{K}$.\footnote{La page écrit
$\mathbb{T}_{K}(A) = \mathbb{T}_{\ell}(A)$, sans $\otimes K$. Elle ne définit
pas $\mathfrak{g}$ ; d'après la suite, où interviennent les Frobenius de
« presque tous » les points, c'est l'algèbre de Lie engendrée par l'image de
Galois pour un modèle de $A$ sur un corps de type fini, et l'invariance sous
$\mathfrak{g}$ est celle des classes de Tate.} Supposons :
\begin{itemize}
  \item[a)] la conjecture de Tate pour les diviseurs : $H^{2}(A,
  \mathbb{Q}_{\ell}(1))^{\mathrm{alg}} = H^{2}(A,
  \mathbb{Q}_{\ell}(1))^{\mathfrak{g}}$ ;
  \item[b)] ou bien les valeurs propres des Frobenius sur $T_{K}(A)$ sont, à des
  racines de l'unité près, des racines carrées de la norme $q$ du point ; ou
  bien $T_{K}(A)$ se décompose en deux morceaux $V_{1} \oplus V_{2}$, propres
  pour presque tous les Frobenius, avec des valeurs propres différentes pour
  presque tous.
\end{itemize}
Alors on doit avoir
\[
D^{2i}(A, \mathbb{Q}_{\ell}(i)) = H^{2i}(A, \mathbb{Q}_{\ell}(i))^{\mathrm{alg}}
= H^{2i}(A, \mathbb{Q}_{\ell}(i))^{\mathfrak{g}} ,
\]
car on a a priori les inclusions de gauche à droite, et l'hypothèse donne
l'égalité des extrêmes.\footnote{La page écrit « $\varepsilon q^{-1/2}$ » ; le
signe de l'exposant dépend de ce qu'on prend le Frobenius arithmétique ou
géométrique. En marge : « par Tate, la première [condition] est un cas
particulier de la seconde ». L'égalité des extrêmes est de l'algèbre
linéaire, que la page ne fait pas : dans le premier cas, après extension finie
le Frobenius agit sur $H^{1}$ par un scalaire, toutes les classes de $H^{2i}(i)$
sont de Tate, et $H^{2i} = \bigwedge^{2i}H^{1}$ est engendré par $H^{2}$ ; dans
le second, avec des valeurs propres $\alpha \neq \beta$, $\alpha\beta = q$,
dont le quotient n'est pas une racine de l'unité, les classes de Tate de
$\bigwedge^{2i}(V_{1} \oplus V_{2})(i)$ forment $\bigwedge^{i}V_{1} \otimes
\bigwedge^{i}V_{2}$, engendré par les produits d'éléments de $V_{1} \otimes V_{2}
\subset H^{2}(1)$. Dans les deux cas, a) fait le reste.}

Ces conditions sont vérifiées si $A$ est une puissance d'une courbe elliptique
définie sur un corps fini $\mathbb{F}_{q}$ : a) se ramène au produit de deux
courbes elliptiques et « a été vérifié par Mumford », et b) est trivial à partir
du cas de la courbe, et stable par puissances cartésiennes. La conjecture de
Tate vaut donc, en tous degrés, pour les puissances d'une courbe elliptique sur
un corps fini.\footnote{L'attribution à Mumford est celle de la page ; nous
ne l'avons pas retracée. Depuis, a) est un théorème pour toutes les variétés
abéliennes sur un corps fini (J. Tate, « Endomorphisms of abelian varieties over
finite fields », \emph{Invent.\ Math.} 2, 1966), et la conjecture de Tate a été
démontrée pour les produits quelconques de courbes elliptiques sur un corps fini
par M. Spiess (1999). La page a d'abord écrit « produit », biffé pour
« puissance ».}

\pagerange{43}{46}
\section*{V. La lettre à Serre du 27 août 1965 (pages 43 à 46)}

Quatre feuillets dactylographiés, datés « 27.8.65 », signés « A.
Grothendieck » : c'est la « lettre (1965) » de l'inventaire.\footnote{Le texte
intégral de la lettre est publié dans la \emph{Correspondance
Grothendieck--Serre} (Société mathématique de France, 2001 ; édition bilingue,
2004). La transcription ne le reproduit pas : elle donne les énoncés des
conjectures sous forme resserrée et un relevé en français du reste. On suit ici
la transcription, sans rien restituer de mémoire. La page 42 est un double pâle
du premier feuillet.} Grothendieck y écrit qu'il retourne les questions de
cycles algébriques, qu'il lui manque toujours le pas décisif, et qu'il veut
énoncer les deux conjectures clés de façon purement géométrique, sans
cohomologie.

Notations : $X$ projective, lisse et connexe de dimension $n$, $Y$ une section
hyperplane lisse, $A^{i}_{\mathrm{alg}}(X)_{\mathbb{Q}} =
A^{i}_{\mathrm{alg}}(X) \otimes \mathbb{Q}$ le groupe des classes de cycles de
codimension $i$ à équivalence algébrique près, tensorisé par $\mathbb{Q}$, et
$\xi \in A^{1}_{\mathrm{alg}}(X)_{\mathbb{Q}}$ la classe de $Y$.\footnote{La
lettre note ce groupe $C^{i}(X)$ ; on garde la notation du dossier entier. La
page 75 emploie la même lettre $C$ pour la même notion.}

\textbf{Conjecture A.} Pour $2i \leqslant n$, la multiplication par
$\xi^{n-2i}$ est un isomorphisme $A^{i}_{\mathrm{alg}}(X)_{\mathbb{Q}}
\xrightarrow{\sim} A^{n-i}_{\mathrm{alg}}(X)_{\mathbb{Q}}$.

Il suffirait, dit-il, de vérifier la surjectivité, et pour $n = 2i + 1$ ; le
premier cas ouvert est $i = 1$, $n = 3$. A donnerait pour ces groupes l'analogue
des théorèmes de Lefschetz comparant $X$ et $Y$ par image directe et inverse,
et impliquerait le théorème de Lefschetz cohomologique fort, $\xi^{n-i} \colon
H^{i} \xrightarrow{\sim} H^{2n-i}$, que la lettre dit non démontré ; seul l'est
le théorème « faible », qui résulte de ce que la dimension cohomologique d'une
variété affine de dimension $n$ est au plus $n$ : pour $U = X - Y$, $H^{i}(X)
\to H^{i}(Y)$ est bijectif pour $i \leqslant n - 2$ et injectif pour $i = n - 1$,
et $H^{i}(Y) \to H^{i+2}(X)$ est bijectif pour $i \geqslant n$ et surjectif pour
$i = n - 1$.\footnote{Le théorème de Lefschetz fort en cohomologie
$\ell$-adique a été démontré par Deligne, « La conjecture de Weil II »
(1980). La majoration de la dimension cohomologique des variétés affines est le
théorème d'Artin de SGA 4 (1963-1964).}

\textbf{Conjecture A'.} Pour $2i \geqslant n - 1$, l'image directe
$A^{i}_{\mathrm{alg}}(Y)_{\mathbb{Q}} \to A^{i+1}_{\mathrm{alg}}(X)_{\mathbb{Q}}$
est surjective : à $\tau$-équivalence près, tout cycle de dimension $j < n/2$
sur $X$ est image d'un cycle à coefficients rationnels de $Y$.

Il dit savoir déduire A de A', sauf erreur, à condition d'avoir A' sur un corps
de base non nécessairement algébriquement clos, ou d'énoncer A pour l'équivalence
cohomologique au prix du caractère purement géométrique. A impliquerait la
« formule de Künneth pour les cycles », donc l'intégralité des coefficients des
polynômes caractéristiques, à des puissances de $p$ près en dénominateur, donc
les conjectures de Weil sauf la question des valeurs absolues des valeurs
propres, qui relève de B ; A lui paraît le minimum pour une définition
utilisable des motifs sur un corps.

\textbf{Conjecture B.} Soit $n = 2m$, et soit
$A^{m}_{\mathrm{alg}}(X)_{\mathbb{Q},\mathrm{prim}}$ le noyau de la
multiplication par $\xi$ de $A^{m}_{\mathrm{alg}}(X)_{\mathbb{Q}}$ dans
$A^{m+1}_{\mathrm{alg}}(X)_{\mathbb{Q}}$. La forme $(-1)^{m}\deg(x x')$ est
définie positive sur ce noyau.

Pour Weil, la positivité suffirait ; la forme définie donnerait en outre que la
$\tau$-équivalence est l'équivalence numérique --- l'équivalence homologique à
coefficients $\mathbb{Q}_{\ell}$ étant coincée entre les deux ---, que les
groupes de classes à $\tau$-équivalence près sont de type fini, que
$A^{m}_{\mathrm{alg}}(X)_{\mathbb{Q}} \otimes \mathbb{Q}_{\ell} \to H^{2m}(X)(m)$
est injectif, et que le rang est majoré par le nombre de Betti $b_{2m}$. A et B
ensemble impliqueraient la semi-simplicité de la catégorie des motifs, et que
l'algèbre des correspondances de $X$ est une $\mathbb{Q}$-algèbre semi-simple de
dimension finie, avec involution et trace. Suivent une généralisation de A à une
correspondance quelconque, la semi-simplicité des opérations de Galois dans les
conjectures de Tate, et une théorie raisonnable des « jacobiennes
intermédiaires », morceaux algébriques des jacobiennes de Weil liés à la
cohomologie de degré impair, par lesquelles la démonstration de B pourrait
passer. A et A' lui semblent préliminaires à B, et il donne une forme
équivalente de A' :

\textbf{Conjecture A''.} Soit $U = X - Y$, affine de dimension $n$. Alors
$A^{i}_{\mathrm{alg}}(U)_{\mathbb{Q}} = 0$ pour $2i > n$.

Ce serait peut-être vrai pour toute variété affine lisse, et, pour toute variété
lisse, tout cycle cohomologiquement équivalent à zéro serait $\tau$-équivalent à
zéro --- conséquence de A et B pour les variétés projectives. Ce qu'il faudrait :
un procédé pour déformer un cycle de dimension pas trop grande et le « pousser
à l'infini ». Il invite Serre à y réfléchir et dit s'y être mis le jour même. En
post-scriptum : $H^{i}(X)$ est la cohomologie $\ell$-adique, éventuellement
tordue, la torsion étant omise des notations.

Ce que l'on sait aujourd'hui tient en trois phrases. Sous la forme de la lettre,
qui porte sur l'équivalence algébrique, A et A' sont en défaut, au « premier cas
ouvert » que la lettre désigne ; la conséquence que la lettre tire de la forme
définie de B, « $\tau$-équivalence $=$ équivalence numérique », est fausse pour
la même raison.\footnote{Griffiths, « On the periods of certain rational
integrals », \emph{Ann.\ of Math.} 90 (1969) : sur une hypersurface quintique
générale $X \subset \mathbb{P}^{4}$ ($n = 3$), la différence de deux droites est
homologiquement équivalente à zéro sans que l'un de ses multiples non nuls soit
algébriquement équivalent à zéro. Donc $\dim
A^{2}_{\mathrm{alg}}(X)_{\mathbb{Q}} \geqslant 2 > 1 = \dim
A^{1}_{\mathrm{alg}}(X)_{\mathbb{Q}}$, et $\xi$ ne peut être surjectif : A
tombe pour $i = 1$, $n = 3$. H. Clemens, « Homological equivalence, modulo
algebraic equivalence, is not finitely generated », \emph{Publ.\ Math.\ IHÉS} 58
(1983) : pour la quintique générale, ce quotient tensorisé par $\mathbb{Q}$ est
de dimension infinie, alors que $A^{1}_{\mathrm{alg}}(Y)_{\mathbb{Q}}$ est de
dimension finie ; A' tombe aussi.} Reformulées pour l'équivalence homologique
--- dans l'exposé de Grothendieck au colloque de Bombay (1968) et dans celui de
Kleiman, « Algebraic cycles and the Weil conjectures » (1968) ---, elles sont
devenues les \emph{conjectures standard} : la conjecture de type Lefschetz, que
l'inverse de l'opérateur de Lefschetz soit algébrique, et la conjecture de type
Hodge, la positivité de B ; elles sont ouvertes en général, la première connue
notamment pour les courbes, les surfaces et les variétés abéliennes, la seconde
en caractéristique nulle, où elle résulte de la théorie de
Hodge.\footnote{Pour les variétés abéliennes, Lieberman (1968) et Kleiman
(1968). En caractéristique $p$, la conjecture de type Hodge est connue pour les
surfaces --- c'est le théorème de l'indice --- et ouverte en général. La
semi-simplicité de la catégorie des motifs numériques est, elle, inconditionnelle
(U. Jannsen, \emph{Invent.\ Math.} 107, 1992).} Enfin les conjectures de Weil ont
été démontrées par Deligne (« La conjecture de Weil I », 1974) sans passer par
elles.

\pagerange{47}{70}
\section*{VI. Sorites sur les variétés abéliennes, Picard, Albanese (pages 47 à
70)}

Le feuillet 47 porte le titre, d'une autre encre que les notes qui suivent. La
rédaction reprend la numérotation à 1 : un « Résumé » (pages 48 à 50), un
paragraphe 1 (pages 51 à 57), un paragraphe 2 (pages 57 à 69), deux questions
(page 70). Elle traite le degré 1 à fond, et c'est la partie démontrée du
dossier.

\pagerange{47}{50}
\subsection*{Le résumé, et les quatre cas (pages 47 à 50)}

On sait attacher à une variété deux variétés abéliennes au moyen de ses
cycles : sa variété de Picard $\mathrm{Pic}^{0}_{X}$ par les diviseurs
algébriquement équivalents à zéro, sa variété d'Albanese $\mathrm{Alb}_{X}$ par
les $0$-cycles. Pour deux variétés abéliennes $A$ et $B$ attachées de l'une ou
l'autre façon à $X$ et à $Y$, un cycle $Z$ de dimension convenable sur $X \times
Y$ donne une application $A(k) \to B(k)$ --- il y a $4 = 2 \times 2$ cas. Il
annonce qu'elle est toujours induite par un homomorphisme de variétés abéliennes
$u_{Z}$, qui ne dépend que de la classe d'équivalence numérique de $Z$ ; que,
réciproquement, avec des cycles à coefficients rationnels, tout homomorphisme
s'obtient ainsi ; et qu'en cohomologie $\ell$-adique, par les identifications
\[
\varphi_{T} \colon H^{1}(\mathrm{Alb}_{T}) \xrightarrow{\ \sim\ } H^{1}(T) ,
\qquad
\psi_{T} \colon H^{2d-1}(T)(d-1) \xrightarrow{\ \sim\ }
H^{1}(\mathrm{Pic}^{0}_{T})
\qquad (d = \dim T) ,
\]
l'homomorphisme cohomologique est un isomorphisme si et seulement si
l'homomorphisme de variétés abéliennes est une isogénie ; alors l'inverse est
défini par une correspondance à coefficients rationnels. C'est le cas, verra-t-on,
du cup-produit par $\xi_{1} \cdots \xi_{n-1}$, pour des classes amples
$\xi_{j}$, de $H^{1}(X)$ dans $H^{2n-1}(X)(n-1)$.

La page 50 dresse le tableau des quatre cas ; on le redessine avec les
notations du dossier.

\emph{Cas I} et \emph{cas II}, transposés l'un de l'autre :

\begin{tikzcd}
H^{1}(Y) \arrow[r] & H^{1}(X) \\
H^{1}(\mathrm{Alb}_{Y}) \arrow[u, "\varphi_{Y}"] \arrow[r] & H^{1}(\mathrm{Alb}_{X}) \arrow[u, "\varphi_{X}"'] \\
\mathrm{Alb}_{Y} & \mathrm{Alb}_{X} \arrow[l]
\end{tikzcd}

\begin{tikzcd}
H^{2m-1}(Y)(m-1) \arrow[r] \arrow[d, "\psi_{Y}"'] & H^{2n-1}(X)(n-1) \arrow[d, "\psi_{X}"] \\
H^{1}(\mathrm{Pic}^{0}_{Y}) \arrow[r] & H^{1}(\mathrm{Pic}^{0}_{X}) \\
\mathrm{Pic}^{0}_{Y} & \mathrm{Pic}^{0}_{X} \arrow[l]
\end{tikzcd}

\emph{Cas III} et \emph{cas IV}, chacun transposé de lui-même :

\begin{tikzcd}
H^{2m-1}(Y)(m-1) \arrow[r] \arrow[d, "\psi_{Y}"'] & H^{1}(X) \\
H^{1}(\mathrm{Pic}^{0}_{Y}) \arrow[r] & H^{1}(\mathrm{Alb}_{X}) \arrow[u, "\varphi_{X}"'] \\
\mathrm{Pic}^{0}_{Y} & \mathrm{Alb}_{X} \arrow[l]
\end{tikzcd}

\begin{tikzcd}
H^{1}(Y) \arrow[r] & H^{2n-1}(X)(n-1) \arrow[d, "\psi_{X}"] \\
H^{1}(\mathrm{Alb}_{Y}) \arrow[u, "\varphi_{Y}"] \arrow[r] & H^{1}(\mathrm{Pic}^{0}_{X}) \\
\mathrm{Alb}_{Y} & \mathrm{Pic}^{0}_{X} \arrow[l]
\end{tikzcd}

Pour passer d'un homomorphisme de variétés abéliennes à un cycle sur $X \times
Y$, il faut, dans le cas I, inverser $\varphi_{Y}$ ; dans le cas II, inverser
$\psi_{X}$ --- « il faut des dénominateurs ? » ; dans le cas III, rien : c'est le
théorème classique des correspondances divisorielles ; dans le cas IV, inverser
$\varphi_{Y}$ et $\psi_{X}$, ce qui introduit des dénominateurs. Pour inverser
$\psi_{X}$ on utilise le théorème de Lefschetz ; pour inverser $\varphi_{X}$, le
théorème de Lefschetz et le cas des variétés abéliennes. C'est le programme des
pages 61 à 69.\footnote{Sur la page, la flèche verticale de droite du cas II
porte une lettre lue $\varphi_{X}$, et le but $H^{2n-1}(X)$ est écrit sans
$(n-1)$ ; c'est $\psi_{X}$, et la torsion est rétablie. Au cas II, la remarque
dit « --- $\varphi_{X}$ », le tiret renvoyant à « inverser » : c'est
$\psi_{X}$ qu'il faut inverser, comme le dit d'ailleurs la dernière ligne de la
page.}

\pagerange{51}{57}
\subsection*{1. Un lemme de théorie des intersections (pages 51 à 57)}

\textbf{Théorème 1.1 (page 51).} Soient $X$, $Y$ projectives, lisses et
connexes sur $k$, de dimensions $n$ et $m$, et $Z$ un cycle de dimension $m$ sur
$X \times Y$. Il existe un unique homomorphisme de variétés abéliennes
\[
u_{Z} \colon \mathrm{Alb}_{Y} \longrightarrow \mathrm{Alb}_{X}
\]
tel que, pour tout $0$-cycle $\mathfrak{A}$ de degré $0$ sur $Y$ pour lequel
$Z_{*}(\mathfrak{A})$ est défini, on ait $\mathrm{alb}_{X}(Z_{*}(\mathfrak{A}))
= u_{Z}(\mathrm{alb}_{Y}(\mathfrak{A}))$. De plus, $u_{Z}$ ne dépend que de la
classe d'équivalence numérique de $Z$.\footnote{Il écrit d'abord
« projectives », puis le biffe ; l'hypothèse est nécessaire aux variétés
d'Albanese comme à l'équivalence numérique, et on la rétablit. Il parle des
variétés d'Albanese « homogènes », lecture incertaine. « De degré $0$ » est
ajouté au-dessus de la ligne.}

\textbf{Corollaire 1.2.} Si $\mathrm{alb}_{Y}(\mathfrak{A}) = 0$ et si
$Z_{*}(\mathfrak{A})$ est défini, alors $\mathrm{alb}_{X}(Z_{*}(\mathfrak{A})) =
0$.\footnote{Le « $0$ » est écrit sur le début d'une fin biffée ; la lecture
est incertaine, mais c'est ce que donne le théorème. C'est l'énoncé qui manque à
la page 28.}

\emph{Unicité.} Il existe un ouvert non vide $V$ de $Y$ tel que $Z_{*}(y)$ soit
défini pour $y \in V(k)$, donc $Z_{*}(\mathfrak{A})$ pour tout $0$-cycle à
support dans $V$ ; et $\mathrm{Alb}_{Y}(k)$ est l'ensemble des sommes de tels
$0$-cycles de degré $0$.

\emph{Existence.} Soit $\alpha_{X} \colon X \to \mathrm{Alb}^{1}_{X}$ le
morphisme canonique dans le torseur de degré $1$, $Z' = (\alpha_{X} \times
\mathrm{id}_{Y})_{*}(Z)$, et $B = \mathrm{Pic}^{0}(\mathrm{Alb}^{1}_{X}) \simeq
\mathrm{Pic}^{0}_{X}$, l'isomorphisme étant le transposé de $\alpha_{X}$. Soit
$\Theta$ la classe de Poincaré sur $B \times \mathrm{Alb}^{1}_{X}$, $\theta$ sa
classe numérique relevée à $B \times \mathrm{Alb}^{1}_{X} \times Y$, et
$\mathfrak{z}''$ la classe numérique de $B \times Z'$. Le cycle
\[
c = \pi_{*}(\mathfrak{z}'' \cdot \theta) , \qquad \pi \colon B \times
\mathrm{Alb}^{1}_{X} \times Y \to B \times Y ,
\]
est une classe numérique de diviseurs sur $B \times Y$, donc une correspondance
divisorielle entre $Y$ et $B$, qui définit un homomorphisme $\mathrm{Alb}_{Y} \to
\mathrm{Pic}^{0}(B)$ ; par la bidualité des variétés abéliennes,
$\mathrm{Pic}^{0}(B) \simeq \mathrm{Alb}_{X}$. C'est $u_{Z}$, et il ne dépend
que de la classe numérique de $Z$.\footnote{Qu'une classe numérique de
diviseurs définisse un homomorphisme bien déterminé repose sur le théorème de
Matsusaka (1957) : pour les diviseurs, équivalence numérique et
$\tau$-équivalence coïncident, et un homomorphisme tué par un entier est nul.
La page n'invoque que « Cartier (bidualité) ». En marge de la définition de
$c'$ : « une référence », entouré. Il écrit $\underline{\mathrm{Pic}}^{00}$ et
$\underline{\mathrm{Pic}}^{0}$ selon les endroits.}

Le Théorème 1.1 résulte alors du \textbf{Corollaire 1.3} (l'homomorphisme
construit satisfait la condition de 1.1), qui résulte du \textbf{Lemme 1.4} :
pour un $0$-cycle $\mathfrak{A}$ de degré $0$ sur $Y$ tel que
$Z_{*}(\mathfrak{A})$ soit défini, l'élément
$\Theta(\mathrm{alb}_{X}(Z_{*}(\mathfrak{A})))$ de $\mathrm{Pic}^{0}(B)$ est
égal à $c'(\mathfrak{A})$. En effet, par la bidualité ensembliste
$\mathrm{Alb}_{X}(k) \to \mathrm{Pic}^{0}(B)$, il suffit de comparer les images
dans $\mathrm{Pic}^{0}(B)$, et $c'(\mathfrak{A}) =
\Theta(u_{Z}(\mathrm{alb}_{Y}(\mathfrak{A})))$ par définition de
$u_{Z}$.\footnote{Le Lemme 1.4 dit « un diviseur de degré $0$ sur $Y$ », là où
il faut le $0$-cycle de degré $0$ du Théorème 1.1. Dans les deux formules de la
page 54, l'exposant ${}^{t}$ devant $Z$ est ajouté après coup.} Et le Lemme 1.4
résulte du \textbf{Lemme 1.5} : il existe un diviseur $D$ dans la classe de
$\Theta$ tel que, posant $C = \pi_{*}(Z'' \cdot (D \times Y))$, on ait
$D(\alpha_{X*}(Z_{*}(\mathfrak{A}))) = C(\mathfrak{A})$ dans $Z^{1}(B)$, les deux
membres étant définis --- « de la théorie d'intersections standard, que nous
nous dispensons ici d'expliciter ».

\textbf{Théorème 1.6 (page 55).} Avec les notations de 1.1, le carré suivant,
à coefficients $\mathbb{Z}_{\ell}$, est commutatif :

\begin{tikzcd}
H^{1}(X) \arrow[r, "{}^{t}Z_{*}"] & H^{1}(Y) \\
H^{1}(\mathrm{Alb}_{X}) \arrow[u, "\varphi_{X}"] \arrow[r, "u_{Z}^{*}"] & H^{1}(\mathrm{Alb}_{Y}) \arrow[u, "\varphi_{Y}"']
\end{tikzcd}

--- ce qu'on vérifie en suivant la construction du point de vue
cohomologique.\footnote{La page note ${}^{t}Z_{*}$ par $a(\gamma^{n}(Z))$,
l'exposant pouvant se lire $m$, et dit $\varphi_{X}^{*}$, $\varphi_{Y}^{*}$ des
« isomorphismes mod torsion » ; ce sont même des isomorphismes sur
$\mathbb{Z}_{\ell}$, comme le montrent les pages 72 et 73.}

\textbf{Corollaire 1.7.} Soient $X$, $Y$ des variétés abéliennes de dimensions
$n$ et $m$. Un homomorphisme $H^{1}(X) \to H^{1}(Y)$ est défini par un cycle de
codimension $n$ sur $X \times Y$ si et seulement s'il est de la forme $u^{*}$
pour un homomorphisme de variétés abéliennes $u \colon Y \to X$ ; $Z$ fixé, $u$
est unique et indépendant de $\ell$.

\textbf{Remarques 1.8.} Le Théorème 1.1 se généralise aux schémas projectifs et
lisses à fibres géométriquement connexes sur une base $S$ quelconque, en partant
d'une classe $Z \in \mathrm{CK}^{n}(X \times_{S} Y)$ --- un « $n$-cocycle au sens
du gradué associé au $K$ » : si $\mathrm{Alb}_{Y}$ existe, et si l'on se donne un
morphisme de $X$ dans un schéma « parabélien » sous un schéma abélien $A$, on
obtient $\mathrm{Alb}_{Y} \to A$ par un élément de $\mathrm{CK}^{1}(\mathrm{Alb}_{Y}
\times_{S} \hat{A})$, donné fibre à fibre par la construction
précédente.\footnote{« Parabélien » et « plus généralement » sont lus sous
réserve ; $\mathrm{CK}$ est le gradué de la $K$-théorie pour la filtration
topologique, qui coïncide avec l'anneau de Chow à torsion près (SGA 6).}

\pagerange{57}{60}
\subsection*{2. Application : l'homomorphisme défini par un $1$-cycle (pages 57
à 60)}

Le titre de la page 57 est « Application : l'algébricité de certains
homomorphismes de VA ou d'espaces de cohomologie $\ell$-adique ».

\textbf{Proposition 2.1.} Soient $X$, $Y$ projectives, lisses et connexes, de
dimensions $n$ et $m$, et $Z$ un $1$-cycle sur $X \times Y$. Il existe un unique
homomorphisme $u \colon \mathrm{Pic}^{0}_{X} \to \mathrm{Alb}_{Y}$ tel que, pour
tout diviseur $D$ algébriquement équivalent à zéro sur $X$ pour lequel
${}^{t}Z_{*}(D) = \mathrm{pr}_{Y*}((D \times Y)\cdot Z)$ est défini,
$u(\mathrm{cl}\,D) = \mathrm{alb}_{Y}({}^{t}Z_{*}(D))$. Il est aussi caractérisé
par la commutativité du carré, à coefficients $\mathbb{Z}_{\ell}$,

\begin{tikzcd}
H^{1}(\mathrm{Alb}_{Y}) \arrow[r, "u^{*}"] \arrow[d, "\varphi_{Y}"'] & H^{1}(\mathrm{Pic}^{0}_{X}) \\
H^{1}(Y) \arrow[r, "Z_{*}"'] & H^{2n-1}(X)(n-1) \arrow[u, "\psi_{X}"']
\end{tikzcd}

et $u$ ne dépend que de la classe numérique de $Z$.\footnote{L'énoncé écrit
$Z$ là où l'hypothèse nomme $Z_{1}$. En marge : « échanger le rôle de $X$ et $Y$
? ». La page note $a(\gamma^{n+m-1}(\mathfrak{z}))$ la flèche du bas, qui va de
$Y$ vers $X$, alors qu'au Théorème 1.6 de la page 55 la même notation allait de
$X$ vers $Y$ ; on écrit $Z_{*}$ et ${}^{t}Z_{*}$ selon la convention du
dossier.}

\emph{Démonstration.} Le composé $\psi_{X} \circ Z_{*} \circ \varphi_{Y}$ est
défini par une correspondance algébrique entre $\mathrm{Alb}_{Y}$ et
$\mathrm{Pic}^{0}_{X}$, indépendante de $\ell$, puisque ses trois facteurs le
sont ; par le Corollaire 1.7 il est de la forme $u^{*}$.\footnote{La page écrit
$\psi_{X}\, a(\gamma(\mathfrak{z}))\, \varphi_{X}^{(1)}$ ; le diagramme appelle
$\varphi_{Y}$.} Pour la formule, on construit $u$ directement. Il suffit de
calculer avec les classes rationnelles, puisque la somme d'Albanese d'un
$0$-cycle rationnellement équivalent à zéro est nulle. Soit $\Theta \in
\mathrm{CH}^{1}(X \times \mathrm{Pic}^{0}_{X})$ la classe de Poincaré, normalisée
par $\Theta(0) = 0$, et
\[
C = \pi_{*}\bigl((1_{Y} \times \Theta) \cdot (Z \times 1)\bigr) \in
\mathrm{CH}^{m}(Y \times \mathrm{Pic}^{0}_{X}) ,
\qquad \pi \colon Y \times X \times \mathrm{Pic}^{0}_{X} \to Y \times
\mathrm{Pic}^{0}_{X} ,
\]
de sorte que ${}^{t}Z_{*}(\Theta_{b}) = C_{*}((b) - (0))$. Le Théorème 1.1
appliqué à $C$ donne $u \colon \mathrm{Pic}^{0}_{X} \to \mathrm{Alb}_{Y}$ avec
$\mathrm{alb}_{Y}({}^{t}Z_{*}(\Theta_{b})) = u(b)$. L'accord avec la cohomologie
se déduit du Théorème 1.6 de la page 55 ; $u$ dépend linéairement de $Z$, et si
$Z$ est numériquement équivalent à zéro, $C$ l'est, donc $u = 0$.

\emph{Variante.} On cherche $\mathrm{Pic}^{0}_{X} \to \mathrm{Alb}_{Y} \simeq
\widehat{\mathrm{Pic}^{0}_{Y}}$, c'est-à-dire une correspondance divisorielle sur
$\mathrm{Pic}^{0}_{X} \times \mathrm{Pic}^{0}_{Y}$ : c'est le composé $\Theta_{Y}
\circ Z \circ {}^{t}\Theta_{X}$.

Un « Corollaire 2.2 » suit, barré de six traits : tout homomorphisme
$\mathrm{Pic}^{0}_{X} \to \mathrm{Alb}_{Y}$ serait défini par un $1$-cycle
entier.\footnote{Il est remplacé par les énoncés à coefficients rationnels des
pages 61 à 69 ; avec des coefficients entiers, c'est le genre de question que
pose la page 70.}

\pagerange{61}{64}
\subsection*{Transposition, et inversion (pages 61 à 64)}

\textbf{Corollaire 2.2 (page 61).} Dans la situation de 2.1, le transposé
$\hat{u} \colon \mathrm{Pic}^{0}_{Y} \to \mathrm{Alb}_{X}$ est l'homomorphisme
défini par ${}^{t}Z$ : pour tout diviseur $D$ algébriquement équivalent à zéro
sur $Y$ tel que $Z_{*}(D)$ soit défini, $\mathrm{alb}_{X}(Z_{*}(D)) =
\hat{u}(\mathrm{cl}\,D)$. Vérification « par exemple par voie cohomologique » ;
en marge, deux fois : « attention aux signes éventuels ! ».\footnote{Le signe
dépend des normalisations de la dualité entre $\mathrm{Alb}$ et
$\mathrm{Pic}^{0}$ --- pour une courbe, selon la convention choisie, le composé
de la polarisation thêta et du tiré en arrière par l'application d'Abel--Jacobi
vaut $1$ ou $-1$. La page le réserve ; nous ne l'avons pas fixé. La page 60
s'achevait sur un autre énoncé sous le même numéro, barré.}

\textbf{Proposition 2.3}, entre crochets et avec en marge « devrait remonter au
n° 1 ». Dans la situation du Théorème 1.1, le transposé $\hat{u}_{Z} \colon
\mathrm{Pic}^{0}_{X} \to \mathrm{Pic}^{0}_{Y}$ vérifie ${}^{t}Z_{*}(D) \sim
\hat{u}_{Z}(\mathrm{cl}\,D)$ pour tout diviseur $D$ algébriquement équivalent à
zéro sur $X$. Ces deux énoncés donnent une autre démonstration de 1.1, par un
homomorphisme de schémas en groupes $\underline{\mathrm{Pic}}_{X/k} \to
\underline{\mathrm{Pic}}_{Y/k}$ ; vérification « par AQT », lecture incertaine.

\textbf{Proposition 2.4.} Dans la situation de 2.1, et $\ell$ fixé,
$Z_{*} \colon H^{1}(Y) \to H^{2n-1}(X)(n-1)$ est un isomorphisme si et seulement
si $u$ est une isogénie. Il existe alors un diviseur $D$ à coefficients
rationnels sur $X \times Y$ tel que ${}^{t}D_{*} \colon H^{2n-1}(X)(n-1) \to
H^{1}(Y)$ soit l'inverse de $Z_{*}$, ce qui caractérise $D$ comme correspondance
divisorielle à coefficients rationnels ; l'homomorphisme
$\mathrm{Alb}_{Y} \to \mathrm{Pic}^{0}_{X}$ qu'il définit est $u^{-1}$, dans
$\mathrm{Hom}(\mathrm{Alb}_{Y}, \mathrm{Pic}^{0}_{X}) \otimes \mathbb{Q}$. On peut
prendre $D$ à coefficients entiers si et seulement si $u$ est un isomorphisme,
et, si $N$ annule le noyau de $u$ comme schéma en groupes, on peut prendre $D =
N^{-1}D_{0}$ avec $D_{0}$ entier.

En effet, dans le carré de 2.1, $\varphi_{Y}$ et $\psi_{X}$ sont des
isomorphismes : $Z_{*}$ en est un si et seulement si $u^{*}$ en est un,
c'est-à-dire si $u$ est une isogénie. L'inverse de $u$ dans
$\mathrm{Hom}\otimes\mathbb{Q}$ est défini, par le théorème des correspondances
divisorielles (cas III), par un diviseur rationnel $D$, dont l'effet
cohomologique est l'inverse de $Z_{*}$. Si $N$ tue $\operatorname{Ker} u$, $N
u^{-1}$ est un vrai homomorphisme ; le meilleur dénominateur est l'exposant de
ce noyau.\footnote{Les lignes de la page sur les dénominateurs sont raturées et
en partie illisibles ; elles parlent d'une forme $N^{-r}D_{0}$ et, dans un ajout
encadré en tête de la page 63, du « meilleur dénominateur possible $b$ étant
donné par $\operatorname{card}\operatorname{Ker} u$ ». Le cardinal est un
dénominateur possible ; le meilleur est l'exposant. La référence « cf … »
est laissée en blanc.}

\textbf{Corollaire 2.5.} Pour $\ell$ premier à la caractéristique, $u$ est une
isogénie de degré premier à $\ell$ si et seulement si le composé
$H^{1}(Y, \mathbb{Z}_{\ell}) \to H^{2n-1}(X, \mathbb{Z}_{\ell})(n-1) \to
H^{2n-1}(X, \mathbb{Z}_{\ell})(n-1)/\mathrm{torsion}$ est bijectif ; il suffit
que la première flèche le soit.\footnote{La page souligne le second $H^{2n-1}$,
ce qu'on lit comme le quotient par la torsion ; celui-ci est $H^{1}(
\mathrm{Pic}^{0}_{X}, \mathbb{Z}_{\ell})$ par la suite exacte de la page 72, et
l'énoncé est alors juste.}

\pagerange{64}{67}
\subsection*{Le théorème de Lefschetz pour $H^{1}$ (pages 64 à 67)}

\textbf{Corollaire 2.6.} Soit $X$ projective, lisse et connexe de dimension $n
\geqslant 2$, et $\xi_{1}, \ldots, \xi_{n-1} \in \mathrm{CH}^{1}(X)$ des classes
amples. Pour tout $\ell$, le cup-produit par $\eta = \xi_{1}\cdots\xi_{n-1}$,
\[
H^{1}(X) \longrightarrow H^{2n-1}(X)(n-1) ,
\]
est un isomorphisme, et son inverse est défini par un élément de
$\mathrm{CH}^{1}(X \times X) \otimes \mathbb{Q}$, unique et indépendant de
$\ell$.

En marge : « Th ? C'est un cas particulier important du “Lefschetz vachement”
donnant “Lefschetz” aux surfaces ».\footnote{« Lefschetz vachement » : le
théorème de Lefschetz fort. C'est ici son cas $H^{1}$, avec en prime
l'algébricité de l'inverse, c'est-à-dire la conjecture standard de type
Lefschetz en degré $1$ ; l'un et l'autre sont démontrés ici par les variétés
abéliennes, à la manière de Weil. Le théorème fort en tout degré, en
cohomologie $\ell$-adique, est de Deligne (1980).} Soit $Z$ le $1$-cycle
$\Delta_{*}(\eta)$ sur $X \times X$ : son effet $Z_{*}$ est le cup-produit par
$\eta$, et par 2.4 il suffit de voir que l'homomorphisme $u \colon
\mathrm{Pic}^{0}_{X} \to \mathrm{Alb}_{X}$ de 2.1 est une isogénie ; mieux, c'est
une polarisation de $\mathrm{Pic}^{0}_{X}$, compte tenu de $\mathrm{Alb}_{X}
\simeq \widehat{\mathrm{Pic}^{0}_{X}}$.\footnote{Une première tentative, par une
chaîne de restrictions à des sections de dimension décroissante
$\mathrm{Pic}^{0}_{X} \to \mathrm{Pic}^{0}_{X_{1}} \to \cdots$, est barrée en
croix.}

Quand $Z$ est l'image d'une courbe lisse projective $v = (v_{1}, v_{2}) \colon C
\to X \times Y$, l'homomorphisme $\mathrm{Pic}^{0}_{X} \to \mathrm{Alb}_{Y}$ de
2.1 est le composé
\[
\mathrm{Pic}^{0}_{X} \xrightarrow{\ v_{1}^{*}\ } \mathrm{Pic}^{0}_{C}
\xrightarrow{\ \text{dualité}\ } \mathrm{Alb}_{C}
\xrightarrow{\ v_{2*}\ } \mathrm{Alb}_{Y} .
\]
Si $X = Y$ et $v_{1} = v_{2} = v$, c'est le tiré en arrière par $v^{*}$ de la
polarisation principale de la jacobienne de $C$, puisque $v_{*}$ est le transposé
de $v^{*}$ ; les conditions suivantes sont donc équivalentes : (i) $u$ est une
polarisation ; (ii) $u$ est une isogénie ; (iii) $v^{*} \colon
\mathrm{Pic}^{0}_{X} \to \mathrm{Pic}^{0}_{C}$ est à noyau fini ; (iv) $v_{*}
\colon \mathrm{Alb}_{C} \to \mathrm{Alb}_{X}$ est surjectif. En cohomologie,
$u^{*}$ est le cup-produit par la classe de $C$, de $H^{1}(\mathrm{Alb}_{X})
\simeq H^{1}(X)$ dans $H^{1}(\mathrm{Pic}^{0}_{X}) \simeq H^{2n-1}(X)(n-1)$.
Or, si $C$ est une intersection générale de diviseurs très amples,
$\mathrm{Alb}_{C} \to \mathrm{Alb}_{X}$ est surjectif --- « fait bien connu ».
Comme on peut remplacer chaque $\xi_{j}$ par un multiple très ample, ce qui
multiplie $\eta$ par un entier non nul, le corollaire est
démontré.\footnote{La réduction aux classes très amples est ajoutée ; la page
suppose d'emblée les diviseurs « très amples ».}

\pagerange{67}{69}
\subsection*{La réciproque (pages 67 à 69)}

\textbf{Proposition 2.7 (page 67).} Soit $D \in \mathrm{CH}^{1}(X \times Y)$,
d'où $u_{D} \colon \mathrm{Alb}_{Y} \to \mathrm{Pic}^{0}_{X}$ et ${}^{t}D_{*}
\colon H^{2n-1}(X)(n-1) \to H^{1}(Y)$. Celui-ci est un isomorphisme si et
seulement si $u_{D}$ est une isogénie ; il existe alors $Z \in \mathrm{CH}_{1}(X
\times Y) \otimes \mathbb{Q}$ tel que $Z_{*}$ soit l'inverse de ${}^{t}D_{*}$,
quel que soit $\ell$, et l'homomorphisme $u_{Z} \colon \mathrm{Pic}^{0}_{X} \to
\mathrm{Alb}_{Y}$ de 2.1 est $u_{D}^{-1}$.

\emph{Démonstration.} On a ${}^{t}D_{*} = \varphi_{Y} \circ u_{D}^{*} \circ
\psi_{X}$ :

\begin{tikzcd}
H^{2n-1}(X)(n-1) \arrow[r, "{}^{t}D_{*}"] \arrow[d, "\psi_{X}"'] & H^{1}(Y) \\
H^{1}(\mathrm{Pic}^{0}_{X}) \arrow[r, "u_{D}^{*}"] & H^{1}(\mathrm{Alb}_{Y}) \arrow[u, "\varphi_{Y}"']
\end{tikzcd}

d'où la première assertion. Pour la seconde, il faut voir que
$\psi_{X}^{-1}$ et $\varphi_{Y}^{-1}$ sont définis par des correspondances à
coefficients rationnels, indépendantes de $\ell$. Choisissons une classe ample
$\xi \in \mathrm{CH}^{1}(X)$, et soit $\lambda$ le cup-produit par $\xi^{n-1}$,
isomorphisme par 2.6. Posons $\mu = \psi_{X} \circ \lambda \circ \varphi_{X}
\colon H^{1}(\mathrm{Alb}_{X}) \to H^{1}(\mathrm{Pic}^{0}_{X})$ :

\begin{tikzcd}
H^{2n-1}(X)(n-1) \arrow[r, "\psi_{X}"] & H^{1}(\mathrm{Pic}^{0}_{X}) \\
H^{1}(X) \arrow[u, "\lambda"] & H^{1}(\mathrm{Alb}_{X}) \arrow[u, "\mu"'] \arrow[l, "\varphi_{X}"]
\end{tikzcd}

C'est l'effet cohomologique de l'isogénie $\mathrm{Pic}^{0}_{X} \to
\mathrm{Alb}_{X}$ de 2.6, de sorte que $\mu^{-1}$ est défini par son inverse dans
$\mathrm{Hom} \otimes \mathbb{Q}$. Donc $\psi_{X}^{-1} = \lambda \circ \varphi_{X}
\circ \mu^{-1}$ et $\varphi_{X}^{-1} = \mu^{-1} \circ \psi_{X} \circ \lambda$ sont
algébriques ; on applique la seconde formule à $Y$.\footnote{Sur la page, la
flèche horizontale du bas a une pointe à chaque bout, la droite surchargée :
le sens n'est pas tranché. On l'oriente de $H^{1}(\mathrm{Alb}_{X})$ vers
$H^{1}(X)$, comme $\varphi_{X}$ partout ailleurs, ce qui est le sens où se lit
la définition de $\mu$. Les références « (cf ) » sont laissées en blanc.}

\pagerange{70}{70}
\subsection*{Deux questions (page 70)}

1) Y a-t-il toujours un cycle \emph{entier} qui inverse $\varphi_{X} \colon
H^{1}(\mathrm{Alb}_{X}) \to H^{1}(X)$ ? Il suffirait que, pour $N$ assez grand,
la fibre générique de $\mathrm{Sym}^{N}(X) \to \mathrm{Alb}^{N}_{X}$ ait un
point rationnel. 2) Y a-t-il toujours un cycle entier qui inverse $\psi_{X}$ ?
Si la réponse à 1) était positive, il suffirait de connaître la réponse à 2)
pour les variétés abéliennes.\footnote{Un point rationnel de la fibre
générique donne une famille de $0$-cycles de degré $N$ paramétrée par
$\mathrm{Alb}_{X}$, de somme l'identité : c'est exactement le cycle $U$ dont la
page 31 disait qu'« a priori on ne sait pas si on peut le trouver ». Pour une
courbe, la réponse est oui : pour $N \geqslant 2g - 1$, $\mathrm{Sym}^{N}C$ est
un fibré projectif sur $\mathrm{Pic}^{N}C$. Le rapprochement est de nous.}

\pagerange{72}{73}
\section*{VII. Deux feuillets : Picard, Albanese et cohomologie entière (pages 72
et 73)}

Deux feuillets isolés, qui donnent sur $\mathbb{Z}_{\ell}$ les identifications
dont le paragraphe 2 s'est servi. Pour $X$ propre sur $k$ et $N$ premier à la
caractéristique, la suite de Kummer donne
\[
H^{1}(X, \mu_{N}) \simeq \mathrm{Pic}(X)[N] , \qquad
0 \to \mathrm{Pic}^{0}(X)[N] \to H^{1}(X, \mu_{N}) \to \mathrm{NS}(X)[N] \to 0 ,
\]
et, à la limite, $H^{1}(X, \mathbb{Z}_{\ell}(1)) \simeq
T_{\ell}(\mathrm{Pic}^{0}_{X})$, le groupe de Néron--Severi ne contribuant pas
puisque $\mathrm{Pic}^{\tau}/\mathrm{Pic}^{0}$ est fini. Si le schéma de Picard a
une composante neutre propre, par exemple si $X$ est normal, on en tire
$\varphi_{X} \colon H^{1}(\mathrm{Alb}_{X}, \mathbb{Z}_{\ell}) \to H^{1}(X,
\mathbb{Z}_{\ell})$, qui est un isomorphisme. Pour $X$ lisse de dimension $n$,
la dualité de Poincaré donne
\[
0 \to \mathrm{NS}(X)[N]^{\vee} \to H^{2n-1}(X, \mathbb{Z}/N)(n-1) \to
H^{1}(\mathrm{Pic}^{0}_{X}, \mathbb{Z}/N) \to 0 ,
\]
et à la limite sur les $\ell^{\nu}$
\[
0 \to \mathrm{NS}(X)\{\ell\}^{\vee} \to H^{2n-1}(X, \mathbb{Z}_{\ell})(n-1)
\xrightarrow{\ \psi_{X}\ } H^{1}(\mathrm{Pic}^{0}_{X}, \mathbb{Z}_{\ell}) \to 0 ,
\]
où $\mathrm{NS}(X)\{\ell\}$ est la composante $\ell$-primaire, finie, du groupe de
Néron--Severi, et la dernière flèche un homomorphisme algébrique.\footnote{Les
torsions à la Tate sont corrigées sur la page, plusieurs fois ; on les donne
sous la forme où les termes s'accordent, $\mathrm{Pic}^{0}(X)[N]^{\vee}$ étant
$\mathrm{Alb}_{X}[N](-1) = H^{1}(\mathrm{Pic}^{0}_{X}, \mathbb{Z}/N)$ par
l'accouplement de Weil.}

Le second feuillet, « Relations entre $\underline{\mathrm{Alb}}^{0}_{X}$,
$\underline{\mathrm{Pic}}^{00}_{X}$ et $H^{1}(X)$, $H^{2n-1}(X)$ », résume : le
pro-$\ell$-groupe fondamental abélianisé de $X$ se surjecte sur celui de
$\mathrm{Alb}_{X}$, avec un noyau fini ; $H^{1}(\mathrm{Alb}_{X}) \to H^{1}(X)$
est un isomorphisme ; $H^{2n-1}(X)(n-1) \to H^{1}(\mathrm{Pic}^{0}_{X}) =
\pi_{1}(\mathrm{Alb}_{X})(-1)$ est surjectif, de noyau fini. Si l'on repère les
variétés abéliennes par leur $H^{1}$, $\mathrm{Alb}_{X}$ correspond à
$H^{1}(X)$ et $\mathrm{Pic}^{0}_{X}$ à $H^{2n-1}(X)(n-1)$.\footnote{La page dit
ces flèches des isomorphismes « à torsion près », et demande entre crochets si
la première est « toujours surjective » : oui. Une isogénie connexe $A' \to
\mathrm{Alb}_{X}$ se tire en arrière en un revêtement connexe de $X$, sans quoi
l'application d'Albanese se relèverait à travers un quotient intermédiaire, en
contradiction avec sa propriété universelle. La réponse est de nous.}

\pagerange{75}{88}
\section*{VIII. Une axiomatique des variétés abéliennes à isogénie près (pages
75 à 88)}

Une suite qu'il numérote « I », cerclé, avec des alinéas numérotés : il
énumère les structures dont est munie la catégorie des variétés abéliennes et
de leurs cycles, pour en faire une axiomatique ; puis il calcule, en termes de
modules de Tate, l'image directe et la classe d'une sous-variété.

\pagerange{75}{77}
\subsection*{Les données (pages 75 et 77)}

(1) La catégorie $\mathcal{VA}$ des variétés abéliennes sur $k$, à isogénie
près ; le foncteur de dualité $A \mapsto \hat{A}$, contravariant ; le foncteur
$A \mapsto C^{*}(A) = A^{*}_{\mathrm{alg}}(A)$ vers les anneaux gradués
commutatifs, à degrés positifs, avec $C^{0} \simeq \mathbb{Z}$ ; la dimension,
additive et invariante par dualité. (2) Le degré $C^{n}(A) \to \mathbb{Z}$,
isomorphisme canonique mais non fonctoriel --- une isogénie de degré $d$
multiplie par $d$ --- et l'isomorphisme fonctoriel $\mathbb{Z} \to C^{0}(A)$ ;
l'image directe des cycles, la formule de projection, qui contient la formule de
transposition. (3) L'« application fondamentale »
\[
C^{1}(A) \longrightarrow \mathrm{Hom}(A, \hat{A})^{\mathrm{sym}} ,
\qquad
C^{1}(A \times B)/\bigl(C^{1}(A) \oplus C^{1}(B)\bigr) \simeq \mathrm{Hom}(A,
\hat{B}) ,
\]
la première se factorisant par $C^{1}(A \times A)$ ; en marge : « on pourrait
définir axiomatiquement $A^{*}$ par cette formule ».\footnote{Il écrit
« $\mathrm{Hom}(A, A^{*})^{\mathrm{alt}}$ » : l'image est formée des
homomorphismes symétriques, dont les formes de Riemann sur les modules de Tate
sont alternées, ce que l'indice rappelle sans doute. Il écrit « $C^{1}(A) \times
C^{1}(B)$ » pour la somme des images réciproques. Après la formule de
projection : « on peut en tirer formellement la loi de Pythagore (?) pour les
cycles », lecture incertaine.}

La réalisation $\ell$-adique : $V_{\ell}(A) = T_{\ell}(A) \otimes \mathbb{Q}_{\ell}$,
covariant, et $H^{1}(A) = \mathrm{Hom}(V_{\ell}(A), \mathbb{Q}_{\ell})$,
contravariant, sont des foncteurs additifs et fidèles, et $H^{*}(A) =
\bigwedge^{*}H^{1}(A)$. La dualité se lit
\[
H^{1}(\hat{A}) \simeq \mathrm{Hom}\bigl(H^{1}(A), \mathbb{Q}_{\ell}(-1)\bigr) ,
\]
fonctoriellement à isomorphisme près, par l'accouplement de Weil. L'application
cycle $\gamma^{2i}_{A} \colon C^{i}(A) \to H^{2i}(A)(i) = \bigwedge^{2i}H^{1}(A)(i)$
est un homomorphisme d'anneaux. Enfin (2) les isomorphismes de trace $\eta_{A}
\colon H^{2n}(A) = \bigwedge^{2n}H^{1}(A) \xrightarrow{\sim} \mathbb{Q}_{\ell}(-n)$ et
$\eta_{\hat{A}}$ se correspondent par la dualité --- « vérifier s'il n'y a pas
un signe ! ».\footnote{Question laissée ouverte sur la page ; la réponse dépend
des conventions pour la puissance extérieure d'un dual et pour l'accouplement de
Weil. Il écrit $C^{*}$ comme foncteur sur $\mathcal{VA}$ vers les groupes
abéliens : covariant pour l'image directe, contravariant pour l'image
réciproque.}

\pagerange{79}{84}
\subsection*{Image directe et classe d'une sous-variété (pages 79 à 84)}

Soit $u \colon A \to A'$ un homomorphisme de variétés abéliennes de dimensions
$n$ et $n'$.\footnote{La page 79 commence par « $u \colon A' \to A$ » et par
« $u^{2n} \colon H^{2n}(A) \to H^{2n}(A')$ » ; toute la suite, jusqu'à la page 84,
ne se lit qu'avec $u \colon A \to A'$ : l'image directe va de $A$ vers $A'$,
est déterminée par $u^{*} \colon H^{1}(A') \to H^{1}(A)$, et $u_{*}(1_{A})$ est
la classe du cycle image dans $A'$. On adopte ce sens partout.} Si $n = n'$,
$u^{*}$ est la multiplication par $d = \deg u$ sur $H^{2n}$, nulle si et
seulement si $u$ n'est pas une isogénie ; de façon équivalente, l'image directe
$u_{*} \colon H^{*}(A) \to H^{*}(A')$, définie par transposition au moyen de la
dualité de Poincaré, vérifie $u_{*}(1_{A}) = d\cdot 1_{A'}$. En général, on
définit par transposition
\[
u_{*} \colon H^{i}(A) \longrightarrow H^{i+2(n'-n)}(A')(n'-n) ,
\]
qui se déduit directement de $u^{*} \colon H^{1}(A') \to H^{1}(A)$ et des
isomorphismes $\eta_{A}$, $\eta_{A'}$. La classe $u_{*}(1_{A}) \in
H^{2(n'-n)}(A')(n'-n)$ est celle du cycle image $u_{*}[A] \in C^{n'-n}(A')$ :
elle est nulle si $u(A)$ est de dimension $< n$, et, si $u(A) = B$ est de
dimension $n$ et $d$ le degré de $A \to B$, elle vaut $d\cdot\gamma(B)$.

En termes de modules de Tate : soit $u_{1} = V_{\ell}(u) \colon V_{\ell}(A) \to
V_{\ell}(A')$. Si $u_{1}$ n'est pas injectif, on trouve $0$ ; s'il l'est, on
trouve un multivecteur de degré $2(n'-n) = \dim\operatorname{Coker} u_{1}$,
c'est-à-dire une forme alternée de degré $2(n'-n)$ sur $V_{\ell}(A')$, nulle sur
$u_{1}(V_{\ell}(A))$, déterminée à un scalaire près, et normalisée ainsi : si $0
\to K \to H^{1}(A') \to H^{1}(A) \to 0$, l'isomorphisme
\[
\bigwedge^{2n'}H^{1}(A') \simeq \bigwedge^{2n}H^{1}(A) \otimes
\bigwedge^{2(n'-n)}K
\]
envoie le produit de la classe fondamentale de $A$ par $u_{*}(1_{A})$ sur $d$
fois la classe fondamentale de $A'$. La partie $\ell$-primaire de $d$ est
l'indice de $u_{1}(T_{\ell}(A))$ dans son saturé.\footnote{La page parle du
« contenu » de $T_{\mathbb{Z}_{\ell}}(A)$ dans $T_{\mathbb{Z}_{\ell}}(A')$ :
sur $\mathbb{Z}_{\ell}$ on ne voit que la partie $\ell$-primaire du degré. Elle
écrit $H^{1}(A'^{*})$ là où la suite exacte demande $H^{1}(A')$, comme le
confirme le $H^{2n'}(A')$ sans astérisque de la ligne suivante, et
$C^{n'-n}(A)$ pour $C^{n'-n}(A')$.}

Par la dualité de Poincaré, $H^{i}(A) \simeq \bigwedge^{2n-i}V_{\ell}(A)(-n)$,
et $u_{*}$ s'identifie à $\bigwedge^{2n-i}u_{1} \otimes \mathbb{Q}_{\ell}(-n)$.
En particulier $u_{*}(1_{A})$ est l'image par $\bigwedge^{2n}u_{1}$ de la classe
fondamentale de $H^{0}(A) = \bigwedge^{2n}V_{\ell}(A)(-n)$ : la connaissance de
$u_{*}(1)$ détermine le sous-espace $u_{1}(V_{\ell}(A))$ de $V_{\ell}(A')$ quand
$u_{1}$ est injectif, comme un multivecteur décomposable détermine le
sous-espace qu'il engendre.\footnote{La page écrit $\bigwedge^{2n-1}u_{1}$ ;
c'est $\bigwedge^{2n-i}$, comme les deux membres l'exigent.}

\pagerange{86}{88}
\subsection*{Homomorphismes et groupe de Néron--Severi (pages 86 et 88)}

Sous le titre « Relations entre hom.\ de V.A.\ et groupes de N.S. », il met en
regard ($n = \dim A$) :

\begin{tikzcd}[column sep=small, nodes={font=\scriptsize}]
C^{1}(A) \arrow[r] \arrow[d] & H^{2}(A)(1) = \bigwedge^{2}H^{1}(A)(1) \arrow[d, "\text{dualité}"] & \\
\mathrm{Hom}(A, \hat{A}) \arrow[r] \arrow[dr, "\text{graphe}"'] & \mathrm{Hom}(H^{1}(\hat{A}), H^{1}(A)) \arrow[dr] & \\
 & C^{n}(A \times \hat{A}) \arrow[r] & H^{2n}(A \times \hat{A})(n)
\end{tikzcd}

où la flèche verticale de droite vient de l'accouplement $H^{1}(A) \times
H^{1}(\hat{A}) \to \mathbb{Q}_{\ell}(-1)$, et l'injection « graphe » n'est « pas
additive ! ». Pour $u \colon A \to \hat{A}$, le graphe est l'image de $\bar{u} =
(\mathrm{id}_{A}, u) \colon A \to A \times \hat{A}$, et $\bar{u}^{*}$ agit sur
$H^{1}(A \times \hat{A}) = H^{1}(A) \oplus H^{1}(\hat{A})$ par
$(\mathrm{id}, u^{*})$.\footnote{La page écrit l'accouplement à valeurs dans
$\mathbb{Q}_{\ell}(1)$ ; avec la convention de la page 77, $H^{1}(\hat{A}) =
\mathrm{Hom}(H^{1}(A), \mathbb{Q}_{\ell}(-1))$, c'est $\mathbb{Q}_{\ell}(-1)$.}

Un homomorphisme $u \colon A \to B$ est connu dès qu'on connaît la classe de
cohomologie de son graphe $\Gamma \subset A \times B$ : par la page 84, elle
détermine $V_{\ell}(\Gamma) \subset V_{\ell}(A \times B)$, donc $\Gamma$. De même,
une sous-variété abélienne est connue par sa classe. L'isomorphisme $C^{1}(A
\times B)/(C^{1}(A) \oplus C^{1}(B)) \simeq \mathrm{Hom}(A, \hat{B})$ se lit en
cohomologie
\[
\frac{\bigwedge^{2}\bigl(H^{1}(A) \oplus H^{1}(B)\bigr)(1)}{\bigwedge^{2}H^{1}(A)(1)
\oplus \bigwedge^{2}H^{1}(B)(1)} \simeq \bigl(H^{1}(A) \otimes H^{1}(B)\bigr)(1)
\simeq \mathrm{Hom}\bigl(V_{\ell}(A), V_{\ell}(\hat{B})\bigr) ,
\]
et il pose la question de prouver que certains éléments de $(H^{1}(A) \otimes
H^{1}(B))(1) \subset H^{2}(A \times B)(1)$ sont algébriques.\footnote{Les mots
qui qualifient ces éléments sont illisibles. Si ce sont les éléments invariants
par Galois, c'est la conjecture de Tate pour les diviseurs sur $A \times B$,
équivalente à $\mathrm{Hom}(A, B) \otimes \mathbb{Q}_{\ell} \simeq
\mathrm{Hom}_{\mathrm{Gal}}(V_{\ell}A, V_{\ell}B)$ : démontrée par Tate (1966) sur
les corps finis, par Zarhin (1975) sur les corps de fonctions de
caractéristique $> 2$, par Faltings (1983) sur les corps de nombres. Sur
$\mathbb{C}$, l'analogue est le théorème de Lefschetz sur les classes de type
$(1,1)$.}

\pagerange{90}{91}
\section*{IX. Le tableau des anneaux de classes de cycles (pages 90 et 91)}

Le feuillet 90 porte seul le titre « Tableau des classes de cycles… » ; le
feuillet 91, « Anneaux de classes de cycles sur une variété projective non
singulière sur un corps $k$ alg.\ clos », est une seule figure, qu'on redessine
avec les notations du dossier.

\begin{tikzcd}[column sep=small, row sep=small, nodes={font=\scriptsize}]
  & & A_{\mathrm{num}} & \\
  & & A_{\tau} \arrow[u, "\overset{?}{\simeq}"'] \arrow[r, "\overset{?}{\simeq}"] & A_{\mathbb{Q}_{\ell}\text{-hom}} \\
  & & A_{\mathrm{alg}} \arrow[u] \arrow[r, "\overset{?}{\simeq}"] & A_{\mathbb{Z}_{\ell}\text{-hom}} \arrow[u] \\
  & & A_{\mathrm{pic}} \arrow[u] & \\
  \mathrm{Gr}_{\mathrm{top}}K \arrow[uuurr] \arrow[uurrr, dashed, "?"'] & & A_{\mathrm{alb}} \arrow[u, "\simeq ?"'] & \\
  \mathrm{Gr}_{\mathrm{alg}}K \arrow[u] & & \mathrm{CH} \arrow[u, "\simeq ?"'] \arrow[ull, "\text{surj}"'] & \\
  & \mathrm{CK} \arrow[ul, "\text{surj}"] \arrow[ur] & &
\end{tikzcd}

Un cadre y entoure la « théorie discrète » --- $A_{\mathrm{num}}$, $A_{\tau}$,
$A_{\mathrm{alg}}$ et les deux équivalences homologiques, à coefficients
$\mathbb{Q}_{\ell}$ et $\mathbb{Z}_{\ell}$ --- ; un autre, la « théorie
continue » --- $A_{\mathrm{pic}}$, $A_{\mathrm{alb}}$ et $\mathrm{CH}$, qu'il
note $A_{\mathrm{lin}}$. À gauche, les gradués de la $K$-théorie et l'anneau
$\mathrm{CK}$, reliés par des « flèches qui sont des isomorphismes mod
torsion ». Des accolades demandent si $A_{\mathrm{num}}$ et $A_{\tau}$ sont des
groupes libres de type fini, si le noyau de $A_{\mathrm{alg}} \to A_{\tau}$ est
fini, si celui de $A_{\mathrm{pic}} \to A_{\mathrm{alg}}$ est une variété
abélienne ; les flèches de la théorie continue sont « toutes surjectives » ; et
une note renvoie à « d'autres équivalences … définies par
Liebermann ».\footnote{L'indice lu $\tau$ est douteux ; les « $G$ » devant $K$
et les deux « surj » sont d'une autre encre. La page ne dit pas quelles
filtrations de $K(X)$ elle note $G_{\mathrm{top}}$ et $G_{\mathrm{alg}}$ ; que
les gradués de la $K$-théorie coïncident avec l'anneau de Chow à torsion près
est le théorème de Riemann--Roch de SGA 6. Une note près de $A_{\mathrm{alg}}
\to A_{\mathbb{Z}_{\ell}\text{-hom}}$ dit « car.\ 0 ou meilleure théorie
intégrale ». Le renvoi est sans doute à D. Lieberman, dont les « Higher Picard
varieties » (1968) comparent ces relations aux jacobiennes intermédiaires de
Weil. L'équivalence de Picard n'est définie nulle part dans le dossier.}

Ce qu'on sait aujourd'hui des points d'interrogation. $A_{\mathrm{num}}$ est
bien libre de type fini. Mais les deux isomorphismes conjecturaux issus de
$A_{\tau}$, vers $A_{\mathrm{num}}$ et vers l'équivalence homologique, sont faux
en général, par l'exemple de Griffiths (voir la page 46) ; $A_{\tau}$ n'est pas de
type fini, même tensorisé par $\mathbb{Q}$ (Clemens) ; et $\mathrm{CH} \to
A_{\mathrm{alb}}$ n'est pas un isomorphisme sur $\mathbb{C}$, par le théorème
de Mumford (voir la page 28). Reste, entre équivalence homologique et
équivalence numérique à coefficients rationnels, la conjecture standard D,
ouverte.\footnote{$A^{i}_{\mathrm{num}}$ est sans torsion, et de rang fini
puisqu'il est quotient de l'image de $A^{i}_{\mathrm{hom}} \otimes \mathbb{Q}$
dans la cohomologie ; comme l'accouplement d'intersection le plonge dans le dual
d'un groupe de type fini, il est libre de type fini. Pour le noyau de
$A_{\mathrm{pic}} \to A_{\mathrm{alg}}$, les Théorèmes 1.4 des pages 28 et 29 et
les pages 72-73 répondent pour les diviseurs et les $0$-cycles ; en
codimension $2$, Murre (1985) a construit un représentant algébrique universel.}

\pagerange{92}{99}
\section*{X. Mumford--Nakai : cycles amples et cônes (pages 92 à 99)}

Le feuillet 92 porte « Mumford-Nakai / Cycles algébriques amples… » ; les
pages 93 à 99 sont un brouillon rapide, dont les formules portent mieux que la
prose.

\pagerange{93}{94}
\subsection*{Le critère (pages 93 et 94)}

\textbf{Théorème de Mumford--Nakai.} Soit $X$ un schéma projectif sur un corps
et $L$ un faisceau inversible sur $X$. Les conditions suivantes sont
équivalentes :
\begin{itemize}
  \item[a)] $L$ est ample ;
  \item[b)] pour tout faisceau cohérent $F$ sur $X$ de support de dimension
  $\geqslant 1$, $\chi(F \otimes L^{\otimes n}) \to +\infty$ quand $n \to
  +\infty$ ;
  \item[b')] de même avec $F = \mathcal{O}_{Z}$, $Z$ sous-schéma fermé intègre
  de dimension $\geqslant 1$ (on peut se borner aux $Z$ normaux) ;
  \item[b'')] pour tout sous-schéma fermé intègre $Z$ de dimension $1$ (resp.\ $\geqslant 2$), il existe $n > 0$ tel que $\chi(L^{\otimes n} \otimes
  \mathcal{O}_{Z}) \geqslant 2$ (resp.\ $\geqslant 1$) ;
  \item[c), c'), c'')] les mêmes, avec $\dim H^{0}$ au lieu de $\chi$.
\end{itemize}
En marge, le critère numérique : pour $X$ irréductible de dimension $d$, $D
\cdot C > 0$ pour toute courbe irréductible $C$ et $D^{d} > 0$, « intersection
calculée dans $GK(X)$ ! ».\footnote{C'est le critère qu'on appelle aujourd'hui
de Nakai--Moishezon (Nakai, 1963 ; Moishezon, 1964), étendu par Kleiman (1966).
Sous sa forme démontrée, il demande $(L^{\dim Z}\cdot Z) > 0$ pour \emph{tout}
sous-schéma fermé intègre $Z$ de dimension $> 0$ --- c'est ce que font les
conditions b') à c'') ---, et pas seulement pour les courbes et pour $X$, comme
la note marginale l'écrit. Une autre note : « le critère implique le critère
d'amplitude de A. Weil ».}

Les implications a) $\Rightarrow$ b) et c) sont connues (théorie du polynôme de
Hilbert, $L$ très ample) ; b) $\Rightarrow$ b') $\Rightarrow$ b'') et c)
$\Rightarrow$ c') $\Rightarrow$ c'') sont triviales. Le reste se fait par
récurrence sur $d = \dim X$. Pour b'') $\Rightarrow$ c''), un \textbf{Lemme 1} :
si les restrictions de $L$ aux sous-schémas intègres de dimension $< d$ sont
amples, $H^{i}(X, L^{\otimes n}) = 0$ pour $i \geqslant 2$ et $n$ grand --- par
restriction aux diviseurs, mais « il faut supposer $X$ projectif » ;
« cf.\ lettre de Mumford ». Pour c'') $\Rightarrow$ a), on se ramène à $X$
intègre ; une section non nulle $f$ de $L^{\otimes n}$ a un lieu des zéros $D$
non vide, puisque $L$ n'est pas trivial, donc de dimension $d - 1$, sur lequel
$L$ est ample par récurrence ; et un \textbf{Lemme 2} : sous les conditions du
Lemme 1, si une puissance de $L$ a une section non nulle, $X \dashrightarrow
\mathbb{P}\,H^{0}(X, L^{\otimes n})$ est partout défini et non constant pour $n$
grand. C'est le schéma de la démonstration de Nakai.

Deux conditions d) et d') sont barrées d'un treillis de traits : l'existence
d'une section non nulle d'une puissance de $L$, jointe à la positivité de $L$
sur les courbes. Il demande si l'on peut y supprimer l'hypothèse de la section,
et si l'on peut supprimer l'hypothèse que $X$ est projectif.\footnote{Pour la
première question, la réponse est non : Mumford a construit, sur une surface
réglée au-dessus d'une courbe de genre $\geqslant 2$, un diviseur de degré $> 0$
sur toute courbe et de carré nul, donc non ample (exemple reproduit par
Hartshorne, \emph{Ample subvarieties of algebraic varieties}, 1970, chap.\ I,
§~10). C'est aussi la question de la page 95.}

\pagerange{95}{95}
\subsection*{Les cônes d'une surface (page 95)}

Sur une surface projective lisse, dans l'espace de Néron--Severi réel, on a le
cône effectif $\mathrm{Eff}$ (les $x$ tels que $nx$ soit la classe d'une courbe
effective pour un $n > 0$), le cône positif $\mathrm{Pos}$ défini par $x^{2}
\geqslant 0$ et $x \cdot a \geqslant 0$ pour une classe ample $a$, et le cône
ample $\mathrm{Amp}$ ; avec leurs intérieurs et leurs adhérences,
\[
\begin{array}{ccccc}
\mathring{\mathrm{Eff}} & \subset & \mathrm{Eff} & \subset & \overline{\mathrm{Eff}} \\
\cup & & & & \cup \\
\mathring{\mathrm{Pos}} & \subset & \mathrm{Pos} & = & \overline{\mathrm{Pos}} \\
\cup & & & & \cup \\
\mathring{\mathrm{Amp}} & = & \mathrm{Amp} & \subset & \overline{\mathrm{Amp}}
\end{array}
\]
L'inclusion encadrée $\mathring{\mathrm{Pos}} \subset \mathrm{Eff}$ --- par
Riemann--Roch, un $x$ de carré $> 0$ a un multiple effectif --- donne
$\mathrm{Pos} \subset \overline{\mathrm{Eff}}$, et, comme le cône positif est
son propre dual pour la forme d'intersection (théorème de l'indice de Hodge),
$\overline{\mathrm{Eff}}^{\vee} = \mathrm{Eff}^{\vee} \subset
\mathrm{Pos}$.\footnote{La page caractérise aussi l'effectivité par « $\ell(nx)
\geqslant 2$ pour $n$ grand » : cette condition dit que $x$ est
\emph{mobile}, non qu'il est effectif --- une courbe exceptionnelle $E$ a
$h^{0}(nE) = 1$ pour tout $n$. La caractérisation juste est $h^{0}(nx)
\geqslant 1$ pour un $n > 0$. Une proposition est citée, au nom illisible.}
Par le théorème de Nakai, $x$ est ample si et seulement si $x \cdot C > 0$ pour
toute courbe effective $C$ et $x^{2} > 0$. D'où
\[
\boxed{\;\mathrm{Amp} = \operatorname{int}\bigl(\mathrm{Eff}^{\vee}\bigr) \subset
\mathring{\mathrm{Pos}} , \qquad \overline{\mathrm{Amp}} = \mathrm{Eff}^{\vee} =
\overline{\mathrm{Eff}}^{\vee} \subset \mathrm{Pos}\;}
\]
: le cône ample est l'intérieur du cône nef, dual du cône effectif, qui est
son adhérence. C'est, en dimension $2$, le théorème de Kleiman (1966).

\emph{Question} : $x \cdot C > 0$ pour toute courbe effective entraîne-t-il $x$
ample ? Par Nakai, il reste à voir $x^{2} \neq 0$ ; on peut supposer $x$ classe
d'un diviseur $D$ avec $D^{2} = 0$ ; alors $\chi(nD) = \chi(\mathcal{O}_{X}) +
n\chi(\mathcal{O}_{D})$, et $h^{0}(nD) \geqslant \chi(nD)$ pour $n$ grand, de sorte
qu'on doit avoir $\chi(\mathcal{O}_{D}) \leqslant 0$ --- et la page s'arrête.
La réponse est non : c'est l'exemple de Mumford cité plus haut, où
$\chi(\mathcal{O}_{D}) < 0$, et l'argument ne pouvait pas aboutir.\footnote{Si
$\chi(\mathcal{O}_{D})$ était $> 0$, un multiple de $D$ serait effectif non nul,
et $D^{2} = \frac{1}{n}\,D\cdot(nD) > 0$ : la contrainte obtenue est juste, mais
elle ne contredit rien.}

\pagerange{97}{99}
\subsection*{En dimension supérieure (pages 97 et 99)}

Pour $X$ lisse de dimension $d$, soit $\mathrm{Eff}_{i}$ le cône des classes
numériques de $i$-cycles effectifs, $\mathrm{Nef}^{i} = \mathrm{Eff}_{i}^{\vee}$
le cône dual dans les classes de codimension $i$, et $\alpha_{i} \colon A^{i}
\to A_{d-i}$ l'identification de Poincaré. Il cherche à dire le critère
numérique par ces cônes --- $D^{d} > 0$ et $D \cdot C > 0$ pour les courbes
effectives, contre « $D \in \mathrm{Nef}^{1}$ et $\alpha_{d-1}D^{d-1} \in
\mathrm{Eff}_{1}$ » --- et pose trois questions.\footnote{En tête de la page 97,
isolé, $H^{2}(L^{\otimes n} \otimes K^{\otimes(2n-1)})$, exposant douteux ; en
bas, $c_{i}(E)$. La page 99 porte des croquis de cônes.}

La première, page 99 : si $D \cdot C \geqslant 0$ pour toute courbe effective,
a-t-on $D^{i}\cdot X_{i} \geqslant 0$ pour tout cycle effectif $X_{i}$ de
dimension $i$, et en particulier $D^{3} \geqslant 0$ sur une variété de dimension
$3$ ? Il obtient ce dernier point, encadré, en développant $(D' + nD)^{3}$ pour
une classe $D'$ auxiliaire. La réponse est oui en général : c'est le théorème
de Kleiman (1966), un diviseur nef est limite de classes amples.

La deuxième, page 97 : a-t-on $\alpha_{d-1}(\mathrm{Nef}^{d-1}) \subset
\mathrm{Eff}_{1}$, c'est-à-dire une classe de courbe positive sur tous les
diviseurs effectifs est-elle effective ? Sous la forme des cônes fermés, la
réponse est oui en caractéristique nulle : le dual du cône pseudo-effectif des
diviseurs est le cône des courbes mobiles (Boucksom, Demailly, Păun et
Peternell, 2004).

La troisième, encadrée page 99 : a-t-on $\mathrm{Nef}^{i} \cdot
\mathrm{Nef}^{d-i} \geqslant 0$ ? C'est équivalent, à adhérence près, à
$\alpha_{i}(\mathrm{Nef}^{i}) \subset \overline{\mathrm{Eff}}_{d-i}$ et à
l'énoncé symétrique, et c'est impliqué par $\mathrm{Nef}^{i} \cdot
\mathrm{Nef}^{j} \subset \mathrm{Nef}^{i+j}$. La réponse est non : Debarre, Ein,
Lazarsfeld et Voisin (2011) ont exhibé, sur certaines variétés abéliennes de
dimension $4$, des classes nef de codimension $2$ qui ne sont pas
pseudo-effectives ; il existe donc deux classes nef de codimension $2$ de
produit négatif.\footnote{Ces réponses sont postérieures ; la page pose les
questions sans les trancher, sauf $D^{3} \geqslant 0$.}

\pagerange{101}{122}
\section*{XI. Conjectures de Lefschetz pour les cycles algébriques (pages 101 à
122)}

Le feuillet 101 porte « Conjectures de Lefschetz pour cycles algébriques »,
suivi d'une accolade qui réunit deux lignes : « Extension aux Jac.\ inter. »
et « Conjectures locales ». Les notations que les pages 102 à 117 emploient ne
sont fixées qu'à la page 116 ; on les traduit dès maintenant : ses
$\mathcal{A}^{i}$ ou $\mathcal{A}^{i}_{r}$ sont $\mathrm{CH}^{i}$, ses $N^{i}$ ou
$\mathcal{A}^{i}_{a}$ sont $A^{i}_{\mathrm{alg}}$, ses $J^{i}$ ou
$\mathcal{A}^{i}_{r}(X)^{a}$ sont $\mathrm{CH}^{i}(X)_{\mathrm{alg}}$ ;
« $\equiv$ » et « mod gpes finis » veulent dire : à des groupes finis près ;
« exc.\ $p$ », à de la $p$-torsion près.

\pagerange{101}{102}
\subsection*{Une surface (page 102)}

Soit $X$ une surface projective lisse et $Y$ une section hyperplane lisse. Alors
$A^{1}_{\mathrm{alg}}(X) \to A^{1}_{\mathrm{alg}}(Y) = \mathbb{Z}$ est surjectif
à un groupe fini près, de noyau de rang $\rho - 1$ ;
$\mathrm{CH}^{1}(X)_{\mathrm{alg}} = \mathrm{Pic}^{0}_{X}(k) \to
\mathrm{Pic}^{0}_{Y}(k)$ est injectif à de la $p$-torsion près ; le noyau $K$ de
$\mathrm{Pic}(X) \to \mathrm{Pic}(Y)$ rencontre $\mathrm{Pic}^{0}(X)$ en un
groupe fini, donc est de type fini, et s'envoie, à un groupe fini près,
injectivement dans la partie primitive de $A^{1}_{\mathrm{alg}}(X)$. « En
général », ajoute-t-il, $K$ est fini, et il présume qu'il en est ainsi si les
courbes des bons pinceaux de sections hyperplanes passant par $Y$ sont
irréductibles.\footnote{L'injectivité à la $p$-torsion près vient de la
surjectivité de $\pi_{1}(Y) \to \pi_{1}(X)$ par Lefschetz, qui rend
$H^{1}(X, \mu_{N}) \to H^{1}(Y, \mu_{N})$ injectif. Le « en général » est
nécessaire : sur la quadrique $\mathbb{P}^{1} \times \mathbb{P}^{1}$ et pour $Y$
une conique lisse, de bidegré $(1,1)$, le faisceau $\mathcal{O}(1,-1)$ est de
degré $0$ sur $Y \simeq \mathbb{P}^{1}$, donc trivial, et $K \simeq \mathbb{Z}$.
L'exemple est de nous. La lecture « partie mobile » de $A^{1}_{\mathrm{alg}}(Y)
/\operatorname{Im}A^{1}_{\mathrm{alg}}(X)$ est douteuse, et une question qui
suit est barrée.}

\pagerange{105}{108}
\subsection*{Les conjectures locales (pages 105 à 108)}

Soit $S = \operatorname{Spec} A$, où $A$ est un anneau local noethérien
strictement hensélien, $f$ un élément de son idéal maximal, $X = S -
\{\text{point fermé}\}$ le spectre épointé, $Y = X \cap V(f)$, et $U = X - Y$ ;
on suppose $X$ et $Y$ réguliers, et l'on note $N$ la dimension de
$A$.\footnote{La page ne dit pas ce qu'est $N$ ; c'est la dimension de $A$,
comme le montrent les exemples $N = 2, 3, 4$. « Noeth.\ » est lu sous
réserve.} Il conjecture :

\begin{itemize}
  \item[(1)] $A^{i}_{\mathrm{alg}}(X) \equiv 0$ pour $i = N/2$ ; et, pour $i
  \geqslant (N+1)/2$, $A^{i}_{\mathrm{alg}}(X) \equiv 0$ et
  $\mathrm{CH}^{i}(X)_{\mathrm{alg}} \equiv 0$, c'est-à-dire $\mathrm{CH}^{i}(X)
  \equiv 0$.
  \item[(2)] $\mathrm{CH}^{i}(X)_{\mathrm{alg}} \to \mathrm{CH}^{i}(Y)_{\mathrm{alg}}$
  est bijectif (exc.\ $p$) pour $i \leqslant (N-2)/2$, injectif pour $i \leqslant
  (N-1)/2$, nul pour $i \geqslant N/2$ --- et alors
  $\mathrm{CH}^{i}(Y)_{\mathrm{alg}} = 0$. $A^{i}_{\mathrm{alg}}(X) \to
  A^{i}_{\mathrm{alg}}(Y)$ est bijectif pour $i \leqslant (N-3)/2$, injectif pour $i
  \leqslant (N-2)/2$, nul à un groupe fini près pour $i \geqslant (N-1)/2$. Donc
  $\mathrm{CH}^{i}(X) \to \mathrm{CH}^{i}(Y)$ est bijectif pour $i \leqslant
  (N-3)/2$ et injectif pour $i \leqslant (N-2)/2$ ; les cas critiques sont $i =
  (N-2)/2$, $(N-1)/2$, $N/2$.
  \item[(3)] Conséquence de (1) et (2) : $\mathrm{CH}^{i}(X) \to
  \mathrm{CH}^{i}(U)$ est toujours un isomorphisme à un groupe fini près, et un
  isomorphisme si $i \leqslant (N-1)/2$ ; $\mathrm{CH}^{i}(Y) \to
  \mathrm{CH}^{i+1}(X)$ est toujours nul à un groupe fini près, et nul si $i
  \leqslant (N-1)/2$.
  \item[(4)] Problème : définir des formes d'intersection $A^{i}_{\mathrm{alg}}(X)
  \times A^{i}_{\mathrm{alg}}(X) \to \mathbb{Z}$, et définir les homomorphismes
  des cas critiques au moyen des cycles évanescents.
\end{itemize}

Pour $i = 1$, c'est le domaine des théorèmes de Lefschetz locaux pour le groupe
de Picard de SGA 2, où Grothendieck avait démontré la conjecture de Samuel : une
intersection complète locale de dimension $\geqslant 4$ est
factorielle.\footnote{SGA 2, exposé XI. Pour $A$ régulier de dimension $N
\geqslant 5$ et $A/f$ à singularité isolée, cela donne $\mathrm{Pic}(X) = 0 =
\mathrm{Pic}(Y)$, en accord avec la bijectivité de (2) pour $i = 1 \leqslant
(N-3)/2$. Le rapprochement est de nous.}

\emph{$N = 2$.} (1) $A^{1}_{\mathrm{alg}}(X)$ est fini : « c'est le théorème de
Mumford--Hodge » ; (2) et (3) sont vides ; le groupe essentiel est
$\mathrm{CH}^{1}(X)_{\mathrm{alg}}$.\footnote{Pour une singularité normale de
surface, la finitude de la partie discrète du groupe des classes vient de ce
que la matrice d'intersection des courbes exceptionnelles d'une résolution est
définie négative (D. Mumford, « The topology of normal singularities of an
algebraic surface », \emph{Publ.\ Math.\ IHÉS} 9, 1961).}

\emph{$N = 3$.} (1) $\mathrm{CH}^{2}(X)$ est fini, en particulier
$\mathrm{CH}^{2}(X)_{\mathrm{alg}} = 0$ ; (2) $\mathrm{CH}^{1}(X)_{\mathrm{alg}}
\to \mathrm{CH}^{1}(Y)_{\mathrm{alg}}$ est injectif (exc.\ $p$). Les groupes
essentiels sont $\mathrm{CH}^{1}(X)_{\mathrm{alg}}$ et $A^{1}_{\mathrm{alg}}(X)$,
les homomorphismes essentiels leurs applications vers
$\mathrm{CH}^{1}(Y)_{\mathrm{alg}}$. Problème : définir un complémentaire
naturel de $\mathrm{CH}^{1}(X)_{\mathrm{alg}}$ dans
$\mathrm{CH}^{1}(Y)_{\mathrm{alg}}$, correspondant aux cycles évanescents, et
donner un sens à la formule $\mathrm{CH}^{1}(X)_{\mathrm{alg}} \simeq
\mathrm{CH}^{1}(Y)_{\mathrm{alg}}^{\pi}$ ; interpréter le théorème de l'indice
de Hodge pour les surfaces projectives comme la négativité d'une forme sur
$A^{1}_{\mathrm{alg}}(X)$.\footnote{Pour $A$ l'anneau local au sommet du cône
affine sur une surface projective $V$, on a $\mathrm{Cl}(A) =
\mathrm{Pic}(V)/\mathbb{Z}h$, et $A^{1}_{\mathrm{alg}}(X)$ est, à des groupes
finis près, le quotient de $\mathrm{NS}(V)$ par la classe hyperplane, identifié à
sa partie primitive : la forme d'intersection y est définie négative, et c'est
exactement le théorème de l'indice. La remarque est de nous.}

\emph{$N = 4$.} (1) $A^{2}_{\mathrm{alg}}(X)$ et $\mathrm{CH}^{3}(X)$ sont finis ;
(2) $\mathrm{CH}^{1}(X)_{\mathrm{alg}} \to \mathrm{CH}^{1}(Y)_{\mathrm{alg}}$ est
bijectif et $A^{1}_{\mathrm{alg}}(X) \to A^{1}_{\mathrm{alg}}(Y)$ injectif (exc.\ $p$), donc $\mathrm{Pic}(X) \to \mathrm{Pic}(Y)$ est injectif.\footnote{La page
conclut « bijectif », ce que les deux lignes précédentes ne donnent pas et que
son propre tableau exclut ($i = 1 > (N-3)/2$). C'est faux : pour $A =
k[[x,y,z,w]]$ et $f = xy - zw$, $\mathrm{Pic}(X) = \mathrm{Cl}(A) = 0$ tandis que
$\mathrm{Pic}(Y) = \mathrm{Cl}(A/f) = \mathbb{Z}$, le point double ordinaire de
dimension $3$ n'étant pas factoriel.} Les groupes essentiels sont
$\mathrm{CH}^{2}(X)_{\mathrm{alg}}$ et $A^{1}_{\mathrm{alg}}(X)$ ; problème :
un complémentaire naturel de $A^{1}_{\mathrm{alg}}(X)$ dans
$A^{1}_{\mathrm{alg}}(Y)$, et un sens à $A^{1}_{\mathrm{alg}}(X) \simeq
A^{1}_{\mathrm{alg}}(Y)^{\pi}$.

En général, pour $N = 2m + 1$, $\mathrm{CH}^{i}(X) \equiv 0$ pour $i \geqslant m
+ 1$ ; pour $N = 2m$, $A^{m}_{\mathrm{alg}}(X) \equiv 0$ --- c'est-à-dire
$\mathrm{CH}^{m}(X)_{\mathrm{alg}} \equiv \mathrm{CH}^{m}(X)$ --- et
$\mathrm{CH}^{i}(X) \equiv 0$ pour $i \geqslant m + 1$.

\pagerange{110}{112}
\subsection*{Les conjectures globales (pages 110 à 112)}

Soit $X$ projective, lisse et connexe de dimension $n$, $Y$ une section
hyperplane lisse, et $E$ l'opérateur de Lefschetz, produit par la classe de
$Y$.

\begin{itemize}
  \item[(1)] \emph{Restriction} $\mathrm{CH}^{i}(X) \to \mathrm{CH}^{i}(Y)$.
  Bijective (exc.\ $p$ ?) si $2i \leqslant n - 2$. Si $2i = n - 1$, injective
  (exc.\ $p$), et bijective sur les sous-groupes des classes $\tau$-équivalentes à
  zéro ; d'où la bijectivité sur les classes algébriquement équivalentes à zéro,
  à la torsion de $A^{i}_{\mathrm{alg}}$ près, et l'injectivité de
  $A^{i}_{\tau}(X) \to A^{i}_{\tau}(Y)$. Si $2i = n$, injective sur les classes
  $\tau$-équivalentes à zéro.\footnote{La page conclut, pour $2i = n - 1$ :
  « $\mathcal{A}^{i}_{a}(X) \to \mathcal{A}^{i}_{a}(Y)$ bijectif », avec la
  notation qu'elle vient de définir pour le groupe à équivalence algébrique
  près. Pris ainsi, c'est faux : pour $X = \mathbb{P}^{3}$ et $Y$ une quadrique
  lisse, $A^{1}_{\mathrm{alg}}(X) = \mathbb{Z}$ et $A^{1}_{\mathrm{alg}}(Y) =
  \mathbb{Z}^{2}$. Ce qui suit de la ligne précédente porte sur le sous-groupe
  des classes algébriquement équivalentes à zéro, et c'est ce qu'on écrit. Un
  premier tableau encadré est biffé.}
  \item[(2)] \emph{Image directe} $\mathrm{CH}_{i}(Y) \to \mathrm{CH}_{i}(X)$, en
  dimension $i$. Bijective (exc.\ $p$ ?) si $2i \leqslant n - 3$. Si $2i = n - 2$,
  surjective (exc.\ $p$), surjective sur les classes $\tau$-équivalentes à zéro,
  et bijective à équivalence algébrique près. Si $2i = n - 1$, surjective à
  $\tau$-équivalence près.
  \item[(3)] \emph{Lefschetz fort.} Pour $2i \leqslant n$, $E^{n-2i} \colon
  A^{i}_{\mathrm{alg}}(X) \to A^{n-i}_{\mathrm{alg}}(X)$ est bijectif à des
  groupes finis près ; pour $2i - 1 \leqslant n$, $E^{n-2i+1} \colon
  \mathrm{CH}^{i}(X)_{\tau} \to \mathrm{CH}^{n-i+1}(X)_{\tau}$ l'est aussi. Donc
  $E^{p} \colon \mathrm{CH}^{i}(X) \to \mathrm{CH}^{i+p}(X)$ est injectif à des
  groupes finis près si $p \leqslant n - 2i$, surjectif si $p \geqslant n - 2i +
  1$.
  \item[(2')] Soit $U = X - Y$, affine et lisse, de sorte que $\mathrm{CH}_{i}(U) =
  \mathrm{CH}_{i}(X)/\operatorname{Im}\mathrm{CH}_{i}(Y)$. Si $2i \leqslant n - 2$,
  $\mathrm{CH}_{i}(U) = 0$ (exc.\ $p$ ?) ; si $2i = n - 1$, $A_{i,\mathrm{alg}}(U)
  = 0$.
  \item[(3')] Soit $E_{X}$ le cône affine épointé sur $X$, quasi-affine et lisse
  de dimension $n + 1$. Si $2i \geqslant n + 2$, $\mathrm{CH}^{i}(E_{X})$ est fini
  et $\mathrm{CH}^{i}(E_{X})_{\mathrm{alg}} = 0$ ; si $2i \geqslant n + 1$,
  $A^{i}_{\mathrm{alg}}(E_{X})$ est fini.
\end{itemize}

(1) et (2) sont la forme, pour les cycles, des théorèmes de Lefschetz faibles ;
(3) est celle du théorème fort ; (2') et (3') en sont des reformulations, par
les suites de localisation.\footnote{En (3'), la page donne à $E$ la dimension
« $2n+1$ » ; c'est $n + 1$. En (2'), elle écrit le quotient avec des indices de
codimension, $\mathcal{A}^{i}_{r}(X)/\operatorname{Im}\mathcal{A}^{i}_{r}(Y)$ ;
en indices de dimension, il est juste. Dans (3), une variante pour l'équivalence
$\tau$ est biffée, et « injectif exc.\ $p$ » est corrigé en « bijectif mod
groupes finis ».} Il note ce qu'on doit alors avoir :
$A^{i}_{\mathrm{alg}}(E_{X})$ est, à des groupes finis près, la partie primitive
$\operatorname{Ker}(E^{n-2i+1} \colon A^{i}_{\mathrm{alg}}(X) \to
A^{n-i+1}_{\mathrm{alg}}(X))$ pour $2i \leqslant n$, et nul pour $2i \geqslant n +
1$ ; $\mathrm{CH}^{i}(E_{X})_{\mathrm{alg}}$ est la partie primitive
$\operatorname{Ker}(E^{n-2i+2} \colon \mathrm{CH}^{i}(X)_{\mathrm{alg}} \to
\mathrm{CH}^{n-i+2}(X)_{\mathrm{alg}})$ de la « variété abélienne
intermédiaire » pour $2i \leqslant n + 1$, et nul pour $2i \geqslant n +
2$.\footnote{La page définit ces parties primitives comme les noyaux de
$L^{n-2i}$ et de $L^{n-2i+1}$ : ce sont les exposants des isomorphismes de
Lefschetz de (3), dont les noyaux seraient finis. Les groupes du cône sont les
conoyaux de $E$, $\mathrm{CH}^{i}(E_{X}) = \mathrm{CH}^{i}(X)/E\cdot
\mathrm{CH}^{i-1}(X)$, que la décomposition primitive identifie aux noyaux de
la puissance suivante ; on corrige les exposants. Une troisième ligne, sur
l'équivalence $\tau$, est barrée.}

\pagerange{113}{113}
\subsection*{Un théorème de dimension cohomologique (page 113)}

\textbf{Théorème 4.} Soit $A$ un anneau local noethérien complet dont les
composantes irréductibles sont d'égale dimension, $f$ un élément de son idéal
maximal, et $U = \operatorname{Spec} A_{f}$. Alors $\operatorname{cd} U
\leqslant \dim A$.

C'est l'énoncé de dimension cohomologique qui porte les conjectures locales, au
sens où la majoration d'Artin porte le Lefschetz faible de la lettre. Le
brouillon, traversé d'un long trait au crayon, esquisse la démonstration : on se
ramène à $A$ intègre, normal, complet, de dimension $n$, et à prouver
$H^{i}(U, \mathbb{Z}/\ell) = 0$ pour $i > n$ ; $A$ est fini sur $A_{0} =
k[[t_{1}, \ldots, t_{n}]]$, avec une extension de corps des fractions qu'on
peut supposer séparable ; on choisit $\delta$ avec $\delta f = t_{n}$, et l'on
descend à l'hensélisé $A'_{0}$ de $k[[t_{1}, \ldots, t_{n-1}]][t_{n}]$, $S$
s'écrivant comme produit fibré d'un revêtement fini normal au-dessus de
$\operatorname{Spec} A'_{0}$. La prose de liaison est en grande partie
illisible.\footnote{Le numéro « 4 » suppose des théorèmes 1 à 3 qui ne sont pas
dans le dossier. L'argument par $k[[t_{1}, \ldots, t_{n}]]$ suppose $A$ d'égale
caractéristique et de corps résiduel séparablement clos, ce que la page ne dit
pas. Des énoncés de ce type sont démontrés dans SGA 4, exposé XIX (Artin), pour
les schémas excellents d'égale caractéristique, et Gabber a établi le théorème
de Lefschetz affine dans sa généralité.}

\pagerange{114}{115}
\subsection*{Vérifications en dimension 2, 3 et 4 (pages 114 et 115)}

Il reprend les conjectures (1), (2), (3) et (2') pour $n = 2$, $3$, $4$, ligne à
ligne, et marque chacune « ok » ou « ? ». Pour $n = 2$ : la restriction
$\mathrm{CH}^{0}(X) \to \mathrm{CH}^{0}(Y)$ et l'injectivité sur
$\mathrm{CH}^{1}(X)_{\tau}$, ok ; $\mathrm{CH}_{0}(Y) \to \mathrm{CH}_{0}(X)$
surjectif, et $A_{0,\mathrm{alg}}(Y) \to A_{0,\mathrm{alg}}(X)$ bijectif (les
deux $\simeq \mathbb{Z}$), ok --- avec en marge : « du moins si équiv.\ alg.\ des $0$-cycles $=$ équiv.\ d'Albanese » ; (3) ok ; (2') $\mathrm{CH}_{0}(U) = 0$,
« ok si équivalence numérique $=$ équiv.\ d'Albanese ». Pour $n = 3$ et $n = 4$,
les lignes sur les $0$-cycles reçoivent la même réserve ; d'autres restent
marquées « ? », avec la légende « question “continue” » (les parties
$\tau$-équivalentes à zéro des groupes de Chow en dimension moyenne) ou
« question “numérique” » ($A^{1}_{\mathrm{alg}} \to A^{2}_{\mathrm{alg}}$ par
$E$ pour $n = 3$ ; $A_{1,\mathrm{alg}}(Y) \to A_{1,\mathrm{alg}}(X)$ pour $n =
4$).\footnote{Pour $n = 2$, la page vérifie la seconde famille de (3) avec les
exposants de la première ($E^{2}$ et $E^{0}$) ; le cas non trivial est $E \colon
\mathrm{CH}^{1}(X)_{\tau} \to \mathrm{CH}^{2}(X)_{\tau} =
\mathrm{CH}_{0}(X)_{\deg 0}$, qui manque. Comme la composée avec
$\mathrm{alb}_{X}$ est l'isogénie $\mathrm{Pic}^{0}_{X} \to \mathrm{Alb}_{X}$ du
Corollaire 2.6 de la page 64, il est bijectif à des groupes finis près si et
seulement si le noyau d'Albanese de $X$ est fini. Pour $n = 3$ et $4$, les
exposants sont les bons.}

Ce que l'on sait aujourd'hui éclaire ces marges. Toutes les réserves portant
sur les $0$-cycles sont l'hypothèse que l'équivalence d'Albanese des $0$-cycles
est l'équivalence rationnelle --- la question de la Remarque 1.5 de la page 28.
Or (2) pour les $0$-cycles, (2') et (3') pour $n = 2$, et la seconde famille de
(3) pour $n = 2$, disent tous que $\mathrm{CH}_{0}(X)$ est, à un groupe fini
près, supporté par une courbe --- la section hyperplane $Y$ ---, et c'est faux
sur $\mathbb{C}$ dès que $p_{g}(X) > 0$.\footnote{Mumford (1969) ; sous la forme
« si $\mathrm{CH}_{0}(X)$ est supporté par une courbe, $p_{g} = 0$ », S. Bloch et
V. Srinivas, « Remarks on correspondences and algebraic cycles », \emph{Amer.\ J. Math.} 105 (1983). Pour (3'), $\mathrm{CH}^{2}(E_{X}) =
\mathrm{CH}_{0}(X)/E\cdot\mathrm{Pic}(X)$, et $E\cdot\mathrm{Pic}(X)$ est contenu
dans l'image de $\mathrm{CH}_{0}(Y)$. Même en dimension $3$, pour $X =
\mathbb{P}^{3}$ et $Y$ une surface quartique, la bijectivité de
$\mathrm{CH}_{0}(Y) \to \mathrm{CH}_{0}(X) = \mathbb{Z}$ demandée par (2) pour $2i
\leqslant n - 3$ est fausse, $Y$ étant une surface K3.} Et les « questions
numériques » tombent sur les exemples de Griffiths et de Clemens : la
quintique générale de $\mathbb{P}^{4}$ met en défaut (3) pour $n = 3$, $i = 1$,
c'est-à-dire la conjecture A de la lettre ; vue comme section hyperplane de
$\mathbb{P}^{4}$ plongé par $\mathcal{O}(5)$, elle met en défaut la bijectivité
de $A_{1,\mathrm{alg}}(Y) \to A_{1,\mathrm{alg}}(X)$ de (2) pour $n = 4$ ; et (2)
pour $2i = n - 1$ est la conjecture A', que Clemens réfute.\footnote{Pour $n =
4$, $X = \mathbb{P}^{4}$, $Y$ la quintique : $A_{1,\mathrm{alg}}(\mathbb{P}^{4}) =
\mathbb{Z}$, alors que le noyau de $A_{1,\mathrm{alg}}(Y) \to \mathbb{Z}$ contient
la différence de deux droites, d'ordre infini (Griffiths, 1969). Ces
rapprochements sont de nous. Pour $i = 1$, la conjecture (1) est, elle, le
théorème de Grothendieck--Lefschetz de SGA 2, exposé XII : $\mathrm{Pic}(X) \to
\mathrm{Pic}(Y)$ est injectif si $n \geqslant 3$ et bijectif si $n \geqslant 4$.}
Grothendieck avait mis ses points d'interrogation exactement là.

\pagerange{116}{117}
\subsection*{Notations, et deux hypothèses de finitude (pages 116 et 117)}

La page 116 fixe les notations : pour $X$ projective et lisse, $H^{*}(X)$ la
cohomologie $\ell$-adique et $P^{*}(X)$ sa partie primitive ;
$\mathrm{CH}^{*}(X)$, $A^{*}_{\mathrm{alg}}(X)$ et
$\mathrm{CH}^{*}(X)_{\mathrm{alg}}$, avec la suite exacte
\[
0 \longrightarrow \mathrm{CH}^{*}(X)_{\mathrm{alg}} \longrightarrow
\mathrm{CH}^{*}(X) \longrightarrow A^{*}_{\mathrm{alg}}(X) \longrightarrow 0 ,
\]
et les parties primitives de $A^{*}_{\mathrm{alg}}(X)$ et de
$\mathrm{CH}^{*}(X)_{\mathrm{alg}}$. « Je ne sais pas s'il y a une définition
raisonnable de la “partie primitive” de $\mathrm{CH}^{*}(X)$. »

La page 117 pose deux hypothèses. 1) Les cycles numériquement équivalents à
zéro, à équivalence rationnelle près, sont décrits par un nombre fini de
« familles algébriques » ; d'où l'égalité de l'équivalence numérique et de la
$\tau$-équivalence, et la finitude du type des $A^{i}_{\mathrm{alg}}(X)$. 2)
$(x, y) - (x, b) - (a, y) + (a, b) \sim 0$. Ensemble, elles entraîneraient que
les $\mathrm{CH}^{i}(X)_{\mathrm{alg}}$ sont des variétés abéliennes, que les
produits $\mathrm{CH}^{i}(X)_{\mathrm{alg}} \times
\mathrm{CH}^{j}(Y)_{\mathrm{alg}} \to \mathrm{CH}^{i+j}(X \times
Y)_{\mathrm{alg}}$ ont une propriété dont le nom est illisible, et que
$\mathrm{CH}_{0}(X)_{\mathrm{alg}} \to \mathrm{Alb}_{X}(k)$ est bijectif ; et 2)
équivaut à la bijectivité de $\mathrm{CH}_{0}(X \times Y)_{\mathrm{alg}} \to
\mathrm{CH}_{0}(X)_{\mathrm{alg}} \times \mathrm{CH}_{0}(Y)_{\mathrm{alg}}$. Les
deux hypothèses sont fausses sur $\mathbb{C}$ : la première par Griffiths et
Clemens --- un ensemble fini de familles ne donne qu'un nombre fini de classes
d'équivalence algébrique ---, la seconde par Mumford, comme on l'a vu à la page
29.\footnote{L'équivalence de 2) avec la bijectivité sur les produits est
correcte : le noyau de $\mathrm{CH}_{0}(X \times Y)_{\mathrm{alg}} \to
\mathrm{CH}_{0}(X)_{\mathrm{alg}} \times \mathrm{CH}_{0}(Y)_{\mathrm{alg}}$ est
engendré par les produits $((x) - (a)) \times ((y) - (b))$, par la même
réduction qu'à la page 29.}

\pagerange{119}{122}
\subsection*{Familles limitées de cycles (pages 119 à 122)}

Soit $X$ projectif et $M$ une partie de $K(X)$. Conditions équivalentes : (i) il
existe des familles limitées de faisceaux localement libres $(F_{i})$, $(G_{i})$
telles que $M$ soit contenu dans l'ensemble des $\gamma(F_{i}) - \gamma(G_{i})$ ;
(ii) il existe des entiers $n$, $N$ et un ensemble fini $P$ de polynômes tels
que tout $\xi \in M$ soit de la forme $\gamma(F) - \gamma(G)$, avec $F$, $G$
quotients localement libres de $\mathcal{O}(-n)^{N}$ de polynômes de Hilbert dans
$P$ ; (iii) il existe des faisceaux localement libres fixes $A$, $B$ tels que les
$\xi + \gamma(A)$ soient représentés par des quotients localement libres de $B$.
(iii) $\Rightarrow$ (ii) est trivial ; (ii) $\Leftrightarrow$ (i) est la
caractérisation des familles limitées ; pour (ii) $\Rightarrow$ (iii), si $G =
\mathcal{O}(-n)^{N}/R$, alors $\xi + \gamma(\mathcal{O}(-n)^{N}) = \gamma(F \oplus
R)$, et la page s'interrompt sur la limitation des $F \oplus R$.\footnote{$R$ est
localement libre, noyau d'une surjection entre localement libres, et parcourt
une famille limitée avec $G$ ; il reste à le présenter comme quotient, ce que
la page ne fait pas. En marge : « les polynômes de Hilbert des $\xi \in M$ sont
en nombre fini ? ».}

Page 120 : $X$ de plus lisse, et $M'$ une partie de $Z^{i}(X)$, les cycles étant
pris à « $K$-équivalence » près. Conditions : (i) $M'$ est limitée, définie par
une famille limitée de cycles ; (ii) $M'$ est numériquement limitée --- un nombre
fini de classes numériques --- et il existe $Z_{0}$ tel que les $Z_{0} + Z$, $Z
\in M'$, soient positifs ; (iii) il existe une partie limitée $M$ de $K(X)$ telle
que tout $Z \in M'$ vérifie $(i-1)!\,Z = c^{i}(\xi)$ pour un $\xi \in M$ ; (iv) il
existe une partie limitée de $K^{(i)}(X)$ dont $M'$ soit contenu dans l'image ;
(v) $M'$ est numériquement limitée. La page 122, au dos d'un feuillet
dactylographié étranger, conclut :
\[
(\mathrm{ii}) \Longrightarrow (\mathrm{i}) \Longleftrightarrow (\mathrm{iv})
\Longrightarrow (\mathrm{iii}) \Longrightarrow (\mathrm{v}) ,
\]
la première implication par la construction des familles limitées de
faisceaux, l'équivalence par la définition de la $K$-équivalence des cycles, la
troisième par « R.R. ».\footnote{C'est le théorème de Riemann--Roch sans
dénominateurs : pour $Z$ intègre de codimension $i$, $c_{j}(\mathcal{O}_{Z}) = 0$
pour $0 < j < i$ et $c_{i}(\mathcal{O}_{Z}) = (-1)^{i-1}(i-1)!\,[Z]$. Le signe
s'absorbe en remplaçant $\xi$ par $-\xi$, ce qui change seulement le signe de
$c_{i}$ quand les classes de Chern intermédiaires sont nulles. La page écrit en
(ii) « $Z \in M$ » pour $Z \in M'$. La première glose se lit mal.} La réciproque
(v) $\Rightarrow$ (i) est fausse en général : une famille limitée n'a qu'un
nombre fini de classes d'équivalence algébrique, et sur la quintique générale
les cycles numériquement équivalents à zéro en ont une infinité (Griffiths). C'est
l'hypothèse 1) de la page 117, sous un autre jour.

\pagerange{125}{127}
\section*{XII. Une chaîne d'énoncés sur les cycles algébriquement équivalents à
zéro (pages 125 à 127)}

La page 125 aligne des énoncés reliés par des doubles flèches, et la page 127
en est le brouillon. Pour des variétés projectives lisses $X$, $Y$ et des points
$a, x \in X$, $b, y \in Y$ :

\begin{itemize}
  \item[1°)] le produit $\mathrm{CH}^{i}(X)_{\mathrm{alg}} \times
  \mathrm{CH}^{j}(X)_{\mathrm{alg}} \to \mathrm{CH}^{i+j}(X)_{\mathrm{alg}}$ est
  nul ; 1° bis) de même pour le produit extérieur vers $X \times Y$ ;
  \item[2°)] le produit extérieur $\mathrm{CH}_{0}(X)_{\mathrm{alg}} \times
  \mathrm{CH}_{0}(Y)_{\mathrm{alg}} \to \mathrm{CH}_{0}(X \times Y)$ est nul ; 2°
  bis) de même avec $X = Y$ ;
  \item[3°)] $(x, y) - (x, b) - (a, y) + (a, b) \sim 0$ ; 3° bis) de même avec $X =
  Y$ ;
  \item[4°)] pour une variété abélienne $A$ et $x, y \in A$, $(x + y) - (x) - (y) +
  (0) \sim 0$, c'est-à-dire que $x \mapsto (x) - (0)$ est un homomorphisme $A(k)
  \to \mathrm{CH}_{0}(A)_{\mathrm{alg}}$ --- « mais surjectif ! » ; 4° bis)
  $\mathrm{CH}_{0}(A)_{\mathrm{alg}} \to A(k)$ est un isomorphisme ;
  \item[5°)] pour tout $X$, $\mathrm{CH}_{0}(X)_{\mathrm{alg}} \to
  \mathrm{Alb}_{X}(k)$, qui est surjectif, est un isomorphisme ;
  \item[6°)] $\mathrm{CH}_{0}(X \times Y)_{\mathrm{alg}} \to
  \mathrm{CH}_{0}(X)_{\mathrm{alg}} \times \mathrm{CH}_{0}(Y)_{\mathrm{alg}}$, qui
  est surjectif, est bijectif ; 6° bis) de même avec $X = Y$.
\end{itemize}

Sur la page, 1° $\Leftrightarrow$ 1° bis $\Rightarrow$ 2°, et les énoncés de 2° à
5° sont reliés par des équivalences, 6° n'étant rattaché à rien. Le brouillon
donne un schéma plus prudent : $[2, 2\,\mathrm{bis}, 3, 3\,\mathrm{bis}]
\Leftrightarrow [4, 4\,\mathrm{bis}] \Rightarrow [5]$, et $5 \Rightarrow [6,
6\,\mathrm{bis}] \Rightarrow [2, \ldots]$, une flèche « th.\ de Weil » et une
réciproque marquée « ? » étant barrées.\footnote{Les exposants de
$\mathcal{A}_{0}(A)$ ont, en 4° et au brouillon, la forme d'un $\tau$ plutôt que
de son $a$ ; pour les $0$-cycles sur une variété connexe, être
$\tau$-équivalent à zéro, algébriquement équivalent à zéro et de degré $0$ sont
la même chose. « Kif-kif » : de même. Le brouillon porte aussi, au crayon et
dans l'autre sens, des identités de groupe et un quotient illisible.}

Les implications qu'on vérifie directement sont celles-ci : 2° équivaut à 3°,
les $0$-cycles algébriquement équivalents à zéro étant engendrés par les $(x) -
(a)$ ; 3° bis entraîne 4°, par image directe le long de l'addition $A \times A
\to A$ ; 4° entraîne 4° bis, puisque la composée $A(k) \to
\mathrm{CH}_{0}(A)_{\mathrm{alg}} \to A(k)$ est l'identité et que la première
flèche est surjective ; et 3° équivaut à 6°, par la réduction de la page 29.
Tous ces énoncés sont faux sur $\mathbb{C}$ : 3°, 6°, et donc 2°, 1° bis et 1°,
pour le produit de deux courbes de genre $\geqslant 1$ ; 4°, 4° bis et 5° pour une
surface abélienne, dont le genre géométrique vaut $1$. C'est le théorème de
Mumford (1969).\footnote{Pour 1° directement : sur $X = C \times C'$,
$\mathrm{pr}_{1}^{*}((x) - (a)) \cdot \mathrm{pr}_{2}^{*}((y) - (b))$ est le cycle
de 3°. Le Théorème 1.6 de la page 29 établit l'analogue de 1° dans
$A^{*}_{\mathrm{alb}}(X)$, où il est vrai.}

\pagerange{130}{131}
\section*{XIII. Le membre générique d'un système linéaire (pages 130 et 131)}

Soit $C \subset X$ un sous-schéma lisse de codimension $j + 1$, lieu de base d'un
système linéaire de dimension $j$ de diviseurs, $f \colon X' \to X$ l'éclaté de
$X$ le long de $C$, et $X' \to S$ le morphisme vers l'espace projectif $S$ du
système, une variété rationnelle. Le diviseur exceptionnel est $P(N) \simeq C
\times S$ ; la fibre $Y(y)$ au-dessus du point générique $y$ de $S$ est le membre
générique du système, et $C_{y} \simeq C$ y est de codimension $j$ :

\begin{tikzcd}[column sep=large]
 & C \times S & \\
C \arrow[d, hook, "\operatorname{codim} j+1"'] & P(N) \arrow[l] \arrow[u, no head, "\simeq" description] \arrow[d, hook, "\alpha"] & C_{y} \arrow[l, hook'] \arrow[d, hook, "\alpha_{y}"] \\
X & X' \arrow[l, "f"] \arrow[d] & Y(y) \arrow[l, hook', "i'_{y}"] \arrow[d] \arrow[ll, bend left=35, "i_{y}"] \\
 & S & y \arrow[l]
\end{tikzcd}

Toute classe de $Y(y)$ se relève en une classe $y'$ sur $X'$, et $f^{*}f_{*}y'
= y' + \alpha_{*}(u)$ par la formule de l'éclatement. Restreignant à la fibre
générique, avec $x = f_{*}(y')$ :
\[
i_{y}^{*}(x) = i'^{*}_{y}(y') + \alpha_{y*}(u_{y}) ,
\]
et, $S$ étant rationnelle, $u_{y}$ provient de $\mathrm{CH}(C)$. D'où
\[
\boxed{\;\mathrm{CH}(Y(y)) = \operatorname{Im}\mathrm{CH}(X) +
\operatorname{Im}\mathrm{CH}(C)\;} ,
\qquad
\mathrm{CH}^{i}(Y(y)) = \operatorname{Im}\mathrm{CH}^{i}(X) +
\operatorname{Im}\mathrm{CH}^{i-j}(C) ,
\]
les groupes de Chow du membre générique étant pris sur le corps $k(S)$ : ils sont
engendrés par les restrictions de $X$ et les images directes du lieu de
base.\footnote{La page écrit $A$ pour les groupes de Chow, et $y'$ à la dernière
ligne là où l'on attend $i'^{*}_{y}(y')$. L'hypothèse que $C$ est le lieu de
base d'un système linéaire, qui fait de $P(N)$ le produit $C \times S$ et de
$X'$ un espace fibré sur $S$, est la lecture la plus simple du diagramme ;
la page écrit seulement « cycle d'incidence » $Z \subset X \times S$ et
« éclaté ». Le mot abrégé en tête de la dernière formule ne se lit pas.}

La page 131, au dos d'une feuille à en-tête de l'IHÉS, porte quelques notes au
crayon : la cuspide $x^{2} = y^{3}$ et sa normalisation $k[t]$, engendrée par
$t^{2}$ et $t^{3}$, le nœud $k[x, y]/(xy)$, et des croquis de points sur des
segments et de deux droites qui se croisent.

\end{document}
